\documentclass[modern,linenumbers]{aastex701}

\usepackage{hyperref}
\usepackage{color}
\usepackage{natbib}
\usepackage{graphicx}
\usepackage{xfrac}
\usepackage{ulem}
\usepackage{xcolor}
\usepackage{amsmath}
\usepackage{mathtools}
\usepackage{afterpage}
\usepackage{longtable}

% color macros for editorial notes
\newcommand{\JM}[1]{{\color{blue}#1}}
\newcommand{\todo}[1]{{\color{red}\textbf{TODO:} #1}}

% frequently used symbols
\newcommand{\fastspecfit}{\texttt{FastSpecFit}}
\newcommand{\fastspec}{\texttt{fastspec}}
\newcommand{\fastphot}{\texttt{fastphot}}
\newcommand{\stackfit}{\texttt{stackfit}}
\newcommand{\kms}{\ensuremath{\mathrm{km~s}^{-1}}}
\newcommand{\flux}{\ensuremath{\mathrm{erg~s}^{-1}~\mathrm{cm}^{-2}}}
\newcommand{\fluxdensity}{\ensuremath{\mathrm{erg~s}^{-1}~\mathrm{cm}^{-2}~\text{\AA}^{-1}}}
\newcommand{\ivardensity}{\ensuremath{\mathrm{erg}^{-2}~\mathrm{s}^{2}~\mathrm{cm}^{4}~\text{\AA}^{2}}}
\newcommand{\ivarflux}{\ensuremath{\mathrm{erg}^{-2}~\mathrm{s}^{2}~\mathrm{cm}^{4}}}

\newcommand{\mgii}{\ion{Mg}{2}}
\newcommand{\oii}{[\ion{O}{2}]}
\newcommand{\oiii}{[\ion{O}{3}]}
\newcommand{\nii}{[\ion{N}{2}]}
\newcommand{\sii}{[\ion{S}{2}]}
\newcommand{\ha}{H\ensuremath{\alpha}}
\newcommand{\hb}{H\ensuremath{\beta}}

\newcommand{\niilam}{[\ion{N}{2}]~$\lambda\lambda6584$}
\newcommand{\oiiilam}{[\ion{O}{3}]~$\lambda5007$}

\newcommand{\mgiidoublet}{\ion{Mg}{2}~$\lambda\lambda2796,2803$}
\newcommand{\oiidoublet}{[\ion{O}{2}]~$\lambda\lambda3726,3729$}
\newcommand{\oiiidoublet}{[\ion{O}{3}]~$\lambda\lambda4959,5007$}
\newcommand{\niidoublet}{[\ion{N}{2}]~$\lambda\lambda6548,84$}
\newcommand{\siidoublet}{[\ion{S}{2}]~$\lambda\lambda6716,31$}

\newcommand{\deltachisq}{{\ensuremath{\Delta\chi^2}}}
\newcommand{\vdisp}{\ensuremath{\sigma_\star}}
\newcommand{\tauv}{\ensuremath{\tau_V}}
\newcommand{\Mstar}{\ensuremath{\mathcal{M}_{\star}}}
\newcommand{\msun}{\ensuremath{\mathcal{M}_{\odot}}}
\newcommand{\logmstar}{\log_{10}(\Mstar/\msun)}

\graphicspath{{./}{figures/}}

\submitjournal{\aj}

\shorttitle{FastSpecFit}
\shortauthors{Moustakas et al.}

\begin{document}

\title{FastSpecFit: Spectrophotometric Modeling of 37 Million
  Extragalactic Targets from the Dark Energy Spectroscopic Instrument
  Second Data Release}

\correspondingauthor{John Moustakas}

\author[0000-0002-2733-4559]{John Moustakas}
\email{jmoustakas@siena.edu}
\affil{Department of Physics \& Astronomy, Siena University,
  515 Loudon Road, Loudonville, NY 12110, USA}

\author[0000-0002-4159-4226]{J.~Buhler}
\email{jbuhler@wustl.edu}
\affiliation{Department of Computer Science and Engineering and
  Genetics, Washington University in St. Louis, One Brookings Drive
  St. Louis, MO 63130, USA}

\author[0000-0002-6867-1244]{D.~Scholte}
\email{dscholte@roe.ac.uk}
\affiliation{Institute for Astronomy, University of Edinburgh, Royal
  Observatory, Blackford Hill, Edinburgh EH9 3HJ, UK}

\author{DESI Collaborators...}
\email{desi@desi.desi}
\affiliation{INSTITUTION, ADDRESS}

\begin{abstract}
\todo{Placeholder abstract.} Optical spectrophotometry opens a wide
window into the physical conditions, dynamical state, and formation
history of galaxies and quasars. \fastspecfit{} is open-source Python
software that extracts this information by jointly modeling each DESI
spectrum and its broadband photometry using physically motivated
stellar continuum and emission-line templates. Measured quantities
include light-weighted stellar ages, velocity dispersions, stellar
masses, star-formation rates, K-corrections, rest-frame absolute
magnitudes and colors, and integrated fluxes and equivalent widths for
nearly 50 emission lines spanning the rest-frame UV through
near-infrared. We describe the algorithms, validate the measurements
using repeat observations and comparisons to complementary codes, and
present the \fastspecfit{} value-added catalog for the DESI Data
Release 2, comprising \todo{N} million extragalactic targets.
\end{abstract}

\keywords{Galaxies --- Surveys --- Catalogs --- Spectrophotometry ---
  Stellar populations --- Emission lines}


%% ============================================================
\section{Introduction}
\label{sec:intro}
%% ============================================================

The Dark Energy Spectroscopic Instrument \citep[DESI;][]{levi13a,
  desi-collaboration16a, desi-collaboration24e} is conducting an
eight-year survey (2021--2029) which will ultimately yield
spectroscopic redshifts for more than 60 million galaxies and quasars
over 17,000~deg$^2$. These data are being made public in a series of
staged releases. The Early Data Release
\citep[EDR;][]{desi-collaboration24d}, comprising 1.7~million unique
spectra from the survey validation (SV) campaign (2020 December--2021
May), was released in 2023 June. Data Release~1
\citep[DR1;][]{desi-collaboration26a}, released in 2025 March,
uniformly reprocessed the SV data and added the first 13 months of the
main-survey sample, for a total of 20.4~million unique spectra. Data
Release~2 (DR2), the subject of this paper, extends the main-survey
data through 2024 April, its first three years, includes a staggering
37 million unique redshifts, and is scheduled for public release in
the first part of 2027. Each successive release has sharpened DESI's
baryon acoustic oscillation distance measurements to sub-percent
precision and currently shows a 2.8--4.2$\sigma$ preference for an
evolving dark energy equation of state over the standard $\Lambda$CDM
concordance cosmological model \citep{adame25b, adame25d, adame25e,
  abdul-karim25a, abdul-karim25b}.

Each spectrum DESI acquires encodes a wealth of astrophysical
information beyond its redshift. The shape and absorption-line content
of the stellar continuum trace the underlying stellar populations,
dust content, and kinematics of a galaxy, while its emission-line
spectrum constrains the star-formation rate, chemical abundances,
ionization state, and gas kinematics, as well as the presence and
strength of an active galactic nucleus (AGN) \citep[and references
  therein]{conroy13a}. Extracting these physical properties at DESI
scale---and combining them with broadband photometry to anchor the
overall spectral energy distribution (SED)---requires a fitting code
that is fast, robust, and automated across an enormous range of
signal-to-noise ratios, redshifts, and galaxy types.

A variety of publicly available codes already perform this kind of
spectrophotometric modeling. Full-spectrum fitting codes---among them
\texttt{pPXF} \citep{cappellari04a, cappellari23a}, \texttt{GandALF}
\citep{sarzi06a, sarzi17a}, \texttt{STARLIGHT}
\citep{cid-fernandes05a}, \texttt{FIREFLY} \citep{wilkinson17a},
\texttt{FIT3D}/\texttt{pyPipe3D} \citep{sanchez16a, sanchez22a},
\texttt{ULySS} \citep{koleva09a}, \texttt{alf} \citep{conroy18a}, and
others---extract the kinematics and properties of the stellar
populations and line-emitting gas; SED-fitting codes such as
\texttt{CIGALE} \citep{boquien19a}, \texttt{Bagpipes}
\citep{carnall18a}, \texttt{Prospector} \citep{johnson21a}, and
\texttt{MAGPHYS} \citep{da-cunha08a} combine photometry (and, in some
cases, spectroscopy) with physically motivated stellar population
synthesis (SPS) models to infer star-formation histories, stellar
masses, and dust properties; and codes like \texttt{PyQSOFit}
\citep{guo18a} are tailored to fitting the UV/optical spectra of
quasars.

Each of these codes makes different trade-offs among flexibility,
physical realism, and computational speed, but none were designed to
be run efficiently on the tens of millions of spectra acquired by
DESI. Existing codes also do not take advantage of the DESI resolution
matrix, which compactly encodes the instrumental resolution and
guarantees uncorrelated uncertainties from pixel-to-pixel
\citep{bolton10a, guy23a}; instead, they generally resample spectra
onto a common wavelength grid or into the rest frame, which violates
one of the core assumptions underpinning maximum-likelihood fitting by
introducing off-diagonal terms into the pixel covariance matrix.

\fastspecfit\footnote{\url{https://fastspecfit.readthedocs.io}} has
been developed to meet both these challenges, scaling efficiently to
tens of millions of spectra while fully exploiting DESI's native,
optimally extracted spectral resolution. Using a forward-modeling
fitting approach, \fastspecfit{} convolves stellar continuum and
emission-line templates with each fiber's resolution matrix and
compares the result directly to the observed spectra, preserving the
uncorrelated pixel statistics that resampling would otherwise destroy
(\S\ref{sec:modeling}). In addition, \fastspecfit{} is engineered for
speed; in its default configuration, it fits a typical DESI spectrum
in under a millisecond on a single CPU core, making it possible to
process the full DESI dataset---and to refit it as the data-reduction
pipeline and SPS models evolve---on modest computing resources. This
speed also makes it practical to estimate parameter uncertainties from
repeated fits to Monte Carlo-resampled data rather than from
computationally expensive posterior sampling
(\S\ref{ssec:uncertainties}), and, via a hybrid message passing
interface (MPI) plus multiprocessing design, to scale from a single
serial process on a laptop to thousands of cores on a supercomputer
(\S\ref{ssec:performance}).

Beyond speed, \fastspecfit{} is designed to be flexible, applying only
the physical constraints the data strictly require and otherwise
letting the data determine the best-fitting model. The continuum model
imposes only minimal priors: non-negative stellar population
coefficients, an energy-conserving dust model, and a freely fitted
stellar velocity dispersion bounded widely enough to accommodate the
full spectral diversity of DESI targets. The emission-line model is
similarly flexible, supporting independent narrow and broad kinematic
components and capturing the broad Balmer-line emission that
distinguishes type~1 AGN whenever the data require it. Unlike codes
that target a specific class or spectral type, \fastspecfit{} applies
this same machinery uniformly to every spectrum, without assuming in
advance whether the target is quiescent, star-forming, a
post-starburst, or hosts an AGN. This flexibility minimizes
classification-driven selection effects and maximizes the discovery
potential of a dataset as large and diverse as DESI, where unusual or
physically ambiguous objects emerge naturally from the fit rather than
requiring special-case handling.

Although this paper presents the first comprehensive description of
\fastspecfit{}'s algorithms, measurements from the earlier EDR and DR1
value-added catalogs (VACs; see Appendix~\ref{sec:appendix_vacs}) have
already enabled a remarkably broad range of extragalactic science by
both DESI collaboration members and the broader scientific
community. We briefly highlight a small, illustrative sample of that
work here. Emission-line fluxes, equivalent widths, and stellar masses
from \fastspecfit{} have been used to identify and characterize
changing-look AGN \citep{guo25a} and to triple the census of dwarf AGN
candidates in the local universe \citep{pucha25a}; to measure
direct-method ($T_e$-based) chemical abundances for nearly 50,000
star-forming galaxies in DESI DR2 \citep{scholte26a}, including two of
the most metal-poor star-forming galaxies ever found (J.~Moustakas
et~al., in preparation), and to characterize the atomic gas sequence
and mass-metallicity relation across five decades in stellar mass
\citep{scholte24a}; to assemble catalogs of post-starburst galaxies
\citep{soto25a} and to infer the physical properties of galaxies in
cosmic voids \citep{rincon25a}; and to measure stellar velocity
dispersions for strong gravitational lens candidates
\citep{huang25a}. The \fastspecfit{} measurements have also been
essential to measurements of the galaxy luminosity function
\citep{moore26a} and stellar mass functions \citep{wang25a}, and to
the modeling of intrinsic alignments \citep{siegel25a} and
photometric-redshift calibration \citep{blanco25a} for the next
generation of weak-lensing surveys.

%\fastspecfit{} operates in two complementary modes. \fastspec{} jointly
%fits DESI spectra and broadband photometry, while \fastphot{} fits
%broadband photometry alone, without any spectroscopic input---useful
%for targets with unreliable spectra or for external, photometry-only
%samples. The two modes share the same template library, dust model,
%and physical output quantities, differing only in which data enter the
%fit and in whether the stellar velocity dispersion is a free parameter
%(\S\ref{sec:modes}).

In this paper we describe the DESI spectroscopic and photometric data
(\S\ref{sec:data}) and the algorithms underlying \fastspecfit{} in
detail, including the spectrophotometric and physical quantities the
code derives (\S\ref{sec:modeling}). We illustrate the complete
fitting sequence applied end-to-end to three representative targets
spanning the range of DESI extragalactic classes in
\S\ref{sec:casestudy}.  We then present the \fastspecfit{} VAC for
DESI DR2, comprising 37 million extragalactic targets, document the
code's computational performance, and validate the measurements
against independent codes (\S\ref{sec:vac}). \todo{Missing
  \S\ref{sec:discussion}.} We conclude with a summary and an outlook
toward future improvements (\S\ref{sec:summary}).

Throughout this paper we adopt a flat $\Lambda$CDM cosmology with
$\Omega_m = 0.315$, consistent with \citet{planck-collaboration20a},
and a \citet{chabrier03a} initial mass function (IMF) spanning
$0.1$--$100$~\msun; we define $h \equiv H_0 / (100~\kms\,{\rm
Mpc}^{-1})$ and express stellar masses and other derived quantities in
units with $h=1$.

%% ============================================================
\section{Spectroscopic and Photometric Data}
\label{sec:data}
%% ============================================================

In this section we describe the DESI spectroscopic
(\S\ref{ssec:spectra}) and broadband photometric
(\S\ref{ssec:photometry}) data that \fastspecfit{} fits, and the
target classes and redshift measurements that define the extragalactic
objects that DESI observes (\S\ref{ssec:targets}).

\subsection{DESI Spectroscopy}
\label{ssec:spectra}

DESI is a highly multiplexed fiber spectrograph mounted on the Mayall
4-m telescope at Kitt Peak National Observatory
\citep{desi-collaboration16b, abareshi22a}.  A wide-field optical
corrector \citep{miller24a} feeds a $3\fdg2$-diameter field of view,
within which 5000 robotically actuated, $1\farcs5$-diameter fibers
\citep{poppett24a} are positioned onto target coordinates and feed ten
three-arm spectrographs covering 3600--9824\,\AA{} in the observed
frame \citep{guy23a}. The focal plane is divided into ten wedge-shaped
``petals'', each read out by all ten spectrographs. A single pointing
of the full 5000-fiber focal plane defines a ``tile'', the basic unit
of the survey's tiling strategy \citep{schlafly23a}. The three cameras
sense blue (B: 3600--5930\,\AA), red (R: 5600--7720\,\AA), and
near-infrared (Z: 7470--9824\,\AA) light, at a spectral resolution of
$\mathcal{R}\equiv\lambda/\Delta\lambda \approx 2000-3500$,
$3400-4800$, and $3800-5200$, respectively \citep{guy23a}.

\citet{guy23a} describe the DESI spectroscopic pipeline in detail (see
also \citealt{desi-collaboration24d, desi-collaboration26a} for
shorter summaries). \fastspecfit{} relies on three per-fiber products
from this pipeline: the one-dimensional spectrum, its corresponding
inverse-variance spectrum, and the resolution matrix $\mathbf{R}$, a
band-diagonal, sparse matrix that quantifies the effective,
wavelength-dependent instrumental resolution. $\mathbf{R}$ is a
natural output of the optimal, ``spectroperfectionism'' algorithm of
\citet{bolton10a}, which DESI uses to extract one-dimensional spectra
with uncorrelated pixels. It is also a key ingredient of
\fastspecfit's forward-modeling philosophy: convolving a perfect,
infinite-resolution model spectrum with $\mathbf{R}$ yields the
equivalent model at DESI's effective spectral resolution (see
\S\ref{sec:modeling}).

The per-fiber spectra are extracted onto an 0.8\,\AA{} linear,
wavelength grid (in vacuum and in the Solar System barycentric frame
of reference), flux-calibrated using a network of F-type standard
stars (typically 10 standard stars per petal), and repeat exposures
are coadded using optimal (inverse-variance) weights. \citet{guy23a}
and \citet{koposov26a} report the residual wavelength-calibration
uncertainty to be $\approx0.025$~\AA{} per sky line, or
$\approx1$\,\kms{} at 7000~\AA. Meanwhile, comparison of the
calibrated spectra against hydrogen-atmosphere white-dwarf models show
a residual flux-calibration scatter of about 2\% over most of the
3600--9800~\AA{} range (rising to $\approx6\%$ shortward of 3700~\AA),
comparable to the $\pm6\%$ scatter between synthesized DESI
spectrophotometry and imaging photometry (\S\ref{ssec:photometry}) for
bright stars \citep{guy23a}. We cross-check these reported calibration
uncertainties for significantly fainter extragalactic targets using
\fastspecfit{} in Appendix~\ref{sec:appendix_diagnostics}.

\subsection{Broadband Photometry}
\label{ssec:photometry}

\fastspecfit{} uses broadband optical and infrared (IR) photometry
from the DESI Legacy Imaging Surveys Data Release 9 (DR9)
\citep[hereafter, Legacy Surveys;][]{dey19a}, which cover the full
DESI footprint using two imaging surveys and three cameras. South of a
declination boundary of approximately $+32\arcdeg$, $grz$ imaging
comes from the DECam Legacy Survey \citep[DECaLS;][]{dey19a}, obtained
with the Dark Energy Camera \citep[DECam;][]{flaugher15a} on the
Blanco 4-m telescope at the Cerro Tololo Inter-American Observatory
(CTIO).  North of this boundary, $g$- and $r$-band imaging come from
the Beijing-Arizona Sky Survey \citep[BASS;][]{zou17a}, which was
conducted with the 90Prime camera \citep{williams04a} on the Bok 2.3-m
telescope, and the $z$-band imaging comes from the Mayall $z$-band
Legacy Survey (MzLS), which used the Mosaic-3 camera \citep{dey16a} on
the Mayall 4-m telescope; both telescopes are located at Kitt Peak
National Observatory (KPNO). Following \citet{myers23a}, we track the
imaging system for each target via the \texttt{PHOTSYS} flag
(\texttt{N} or \texttt{S}) to assign the appropriate filter
curves.\footnote{\url{https://speclite.readthedocs.io/en/latest/}}

The Legacy Survey catalogs are generated using the Tractor
forward-modeling photometry pipeline \citep{lang16b, lang16a}, which
fits a common two-dimensional image model simultaneously across all
available $grz$ imaging to derive the integrated flux for each
source. A force-photometry technique is used to derive fluxes in the
four WISE bands---$W1$ (3.4\,\micron), $W2$ (4.6\,\micron), $W3$
(12\,\micron), and $W4$ (22\,\micron) \citep{wright10a}---based on the
unWISE full-depth coadds produced by \citet{lang14a,
  meisner21a}. \fastspecfit{} fits this total photometry, in all seven
$grz$+WISE bands, together with the DESI spectrum in a joint
spectrophotometric fit (\S\ref{sec:modeling}); in any band where a
source is undetected, the reported $1\sigma$ flux uncertainty is used
as an upper limit.

\subsection{DESI Target Classes and Redshifts}
\label{ssec:targets}

This section briefly summarizes the DESI survey and program structure,
the extragalactic target classes, and the redshift-determination
pipeline needed to interpret the \fastspecfit{} VACs; for a much more
detailed discussion of all these topics see
\citet{desi-collaboration24e, desi-collaboration24d,
  desi-collaboration26a}.

DESI observations are organized by survey and program
\citep{schlafly23a, myers23a}. \todo{Need to reconcile the following
  statements with the eight-year program mentioned in the
  introduction. Mention the DESI extension.} The primary five-year
program is the main survey, which began on 2021 May 14 and was
preceded by several SV phases, most notably the One-Percent Survey
(SV3), which covered roughly 1\% of the final DESI footprint to high
spectroscopic completeness \citep{desi-collaboration24e,
  desi-collaboration24d}. Within each survey, observations are
assigned to the ``bright'' program (moonlight or reduced
transparency), the ``dark'' program (nominal dark-sky conditions), or
the ``backup'' program \citep[twilight and poor-weather
  time;][]{dey26a}, on the basis of sky conditions. A small fraction
of observations that do not fit any of these three categories are
assigned to an ``other'' category \citep{desi-collaboration26a}.

DESI photometrically selects primary extragalactic targets from the
Legacy Surveys \citep{dey19a} and organizes them into four classes
spanning complementary redshift ranges \citep{myers23a}. The bright
galaxy survey \citep[BGS, $0 < z \lesssim 0.6$;][]{hahn23a} targets
galaxies to $r \lesssim 20.2$~mag during bright conditions; the
luminous red galaxy \citep[LRG, $0.4 < z \lesssim 1.1$;][]{zhou23a},
emission-line galaxy \citep[ELG, $0.6 < z \lesssim
  1.6$;][]{raichoor23a}, and quasar \citep[QSO, $0.9 < z \lesssim
  4$;][]{chaussidon23a} samples are all observed under dark
conditions; and finally secondary or tertiary (special) targets are
observed as part of a wide range of bespoke science programs and can
be assigned to any program \citep{myers23a, desi-collaboration26a}.

Spectroscopic redshifts and spectral classifications are determined by
Redrock (S.~Bailey et~al., in preparation; \citealt{anand24a}), a
principal component analysis (PCA) template-fitting code that
independently fits templates for stars, galaxies, and quasars over a
type-dependent grid of trial redshifts (S.~Bailey et al., in
preparation; \citealt{guy23a, anand24a}). The two primary
redshift-quality diagnostics reported by Redrock are \texttt{ZWARN}, a
per-object warning bitmask where \texttt{ZWARN=0} signals no known
hardware, data-quality, or redshift-fitting problems, and
\texttt{DELTACHI2}, the difference in $\chi^2$ between the best and
second-best fits across all spectral classes. \todo{Should probably
  mention \texttt{COADD\_FIBERSTATUS}, if we are going to use it in
  \S\ref{ssec:cuts} to select ``good'' redshifts.} Three after-burner
processing steps follow Redrock: an emission-line fitter
(\texttt{emlinefit}; \citealt{raichoor23a}) used to validate ELG
redshifts; a \mgiidoublet{} doublet fitter \citep{chaussidon23a}; and
the \texttt{QuasarNET} neural-network QSO classifier \citep{busca18a,
  green25a}. The latter two after-burners are essential for correctly
classifying and measuring accurate redshifts for QSOs (see
\S\ref{ssec:catalog}).

%% ============================================================
\section{Spectrophotometric Modeling}
\label{sec:modeling}
%% ============================================================

This section describes the technical, algorithmic structure of
\fastspec, the default spectrophotometric fitting mode within the
larger \fastspecfit{} software package. Readers interested in worked
examples of \fastspec{} applied to three illustrative DESI targets may
safely skip to \S\ref{sec:casestudy}, while those interested only in
the resulting DR2 VAC may skip to \S\ref{sec:vac}.

Let $F(\lambda;\theta)$ be the observed-frame ultraviolet (UV) through
IR SED of a galaxy, where $\theta$ denotes the physical
properties---stellar mass, star-formation history, dust content,
kinematics, redshift, and so on---that determine it.
%In practice, DESI fiber spectroscopy captures only $f(\lambda;\theta)
%\equiv F(\lambda;\theta)/\eta$, where $\eta$ is the
%aperture-correction factor relating the flux enclosed within the fiber
%to the galaxy's total flux (generally $\eta \gtrsim 1$, although in
%practice it can scatter below unity; see \S\ref{sssec:apercorr}). We
%rescale the observed fiber spectrum by $\eta$ before fitting, so that
%our data---the spectral pixels $d_i$ and broadband photometry $D_k$
%(\S\ref{sssec:fitting})---are on the same total-flux scale as $F$.
Our overarching goal is to find the model
\begin{equation}
  M(\lambda;\theta) = C(\lambda;\theta_C) + E(\lambda;\theta_E),
  \label{eq:totmodel}
\end{equation}
the sum of a stellar continuum model $C$ (\S\ref{ssec:continuum}) and
an emission-line model $E$ (\S\ref{ssec:emlines}), that best
reproduces $F$ in a maximum-likelihood sense, where $\theta \equiv
\theta_C \cup \theta_E$ is the union of the continuum and
emission-line parameter sets.

%Table~\ref{table:notation} summarizes this and the additional notation
%introduced throughout the section. \todo{I am unlikely to keep this
%  table in its current form. I'll either remove it or shorten it
%  considerably.}

%\startlongtable
\begin{deluxetable*}{lll}
\tablecaption{Notation Used in Section~\ref{sec:modeling}\label{table:notation}}
\tablewidth{0pt}
\tabletypesize{\scriptsize}
\tablehead{
  \colhead{Symbol} &
  \colhead{Description} &
  \colhead{First used}
}
\startdata
\cutinhead{Global Spectrophotometric Model}
$F(\lambda;\theta)$ & Observed-frame UV-through-IR spectral energy distribution & \S\ref{sec:modeling} \\
$\theta$, $\theta_C$, $\theta_E$ & Full / continuum / emission-line parameter sets & \S\ref{sec:modeling} \\
$M(\lambda;\theta)$ & Total spectrophotometric model, $M=C+E$ & \S\ref{sec:modeling} \\
$C(\lambda;\theta_C)$ & Stellar continuum model & \S\ref{sec:modeling} \\
$E(\lambda;\theta_E)$ & Emission-line model & \S\ref{sec:modeling} \\
\cutinhead{SPS Templates and Dust Emission}
$S_n(\lambda)$ & Rest-frame SPS template $n$ & \S\ref{ssec:templates} \\
$S_\mathrm{dust}(\lambda)$ & Infrared dust-emission template & \S\ref{ssec:templates} \\
$c_n$, $\mathbf{c}$ & NNLS coefficient of template $n$; coefficient vector & \S\ref{ssec:templates} \\
$q_\mathrm{PAH}$ & PAH mass fraction & \S\ref{ssec:templates} \\
$U_\mathrm{min}$ & Minimum starlight intensity heating the diffuse ISM & \S\ref{ssec:templates} \\
$\gamma$ & PAH photodissociation-region mass fraction & \S\ref{ssec:templates} \\
$\mathcal{R}_\mathrm{C3K}$ & Resolving power of the C3K template library & \S\ref{ssec:templates} \\
$N$ & Number of SPS age-bin templates ($=5$) & \S\ref{ssec:templates} \\
$\Delta v$ & SPS template pixel size, $=50$~\kms{} (constant-velocity grid) & \S\ref{ssec:templates} \\
\cutinhead{Stellar Continuum Fit}
$S_n^\mathrm{att}$, $S_n^\mathrm{tot}$, $S_n^\mathrm{obs}$ & Attenuated / attenuated$+$reradiated / observed-frame template $n$ & \S\ref{sssec:fitting} \\
$k(\lambda)$ & Dust attenuation curve & \S\ref{sssec:fitting} \\
$L_{\mathrm{abs},n}$ & Luminosity absorbed from template $n$ & \S\ref{sssec:fitting} \\
$\tauv$ & $V$-band dust optical depth & \S\ref{sssec:fitting} \\
$T_\mathrm{IGM}(\lambda,z)$ & IGM transmission & \S\ref{sssec:fitting} \\
$\mathcal{F}_0$ & Dimensionless flux-unit scale factor & \S\ref{sssec:fitting} \\
$\mathcal{M}_0$ & Fixed mass normalization, $10^{10}\,\msun$ & \S\ref{sssec:fitting} \\
$z$ & Redshift & \S\ref{sssec:fitting} \\
$D_L$ & Luminosity distance & \S\ref{sssec:fitting} \\
$n_\mathrm{pix}$, $n_\mathrm{band}$ & Number of spectral pixels in the DESI spectrum / photometric bands (=7) & \S\ref{sssec:fitting} \\
$d$, $d_i$ & DESI spectrum array; flux density at pixel $i$ (reused as $d_j$ in \S\ref{sssec:emline_model}) & \S\ref{sssec:fitting} \\
$I$, $I_i$ & DESI inverse-variance array; value at spectral pixel $i$ & \S\ref{sssec:fitting} \\
$m$, $m_i$ & Binary emission-line mask array (\S\ref{ssec:patches}); value at pixel $i$ ($1$=retained, $0$=masked) & \S\ref{sssec:fitting} \\
$D$, $D_k$ & Tractor total-flux photometry array; value at band $k$ (all 7 bands in the fit; restricted to $grz$ in \S\ref{sssec:apercorr}) & \S\ref{sssec:fitting} \\
$I_D$, $I_{D,k}$ & Photometric inverse-variance array; value at band $k$ & \S\ref{sssec:fitting} \\
$C_i$, $C_k$ & Continuum model evaluated at spectral pixel $i$ / photometric band $k$ & \S\ref{sssec:fitting} \\
$\boldsymbol{\Psi}$ & Continuum design matrix (NNLS) & \S\ref{sssec:fitting} \\
$\mathbf{C}$ & Continuum model evaluated at all pixels/bands, $\mathbf{C}=\boldsymbol{\Psi}\mathbf{c}$ & \S\ref{sssec:fitting} \\
$\mathbf{b}$ & Stacked, weighted spectroscopic$+$photometric data vector (NNLS right-hand side) & \S\ref{sssec:fitting} \\
$\chi^2$ & Fit statistic (reused for the continuum, line, and kinematic-model-selection fits) & \S\ref{sssec:fitting} \\
\cutinhead{Velocity Dispersion and Aperture Correction}
$\vdisp$, $\vdisp^\mathrm{kern}$ & Stellar velocity dispersion; fitted Gaussian kernel width & \S\ref{sssec:vdisp} \\
$\sigma_\mathrm{C3K}$ & Intrinsic velocity resolution of the SPS templates & \S\ref{sssec:vdisp} \\
$\Mstar$ & Stellar mass (preliminary estimate used for the $\vdisp$ fallback) & \S\ref{sssec:vdisp} \\
$\eta_k$, $\bar\eta$ & Per-band / scalar aperture-correction factor & \S\ref{sssec:apercorr} \\
$C_k^\mathrm{fib}$ & Fiber-normalized continuum model, synthesized in band $k$ & \S\ref{sssec:apercorr} \\
\cutinhead{Emission-Line Profile and Kinematics}
$\phi(\lambda)$ & Intrinsic log-Gaussian emission-line profile & \S\ref{sssec:emline_model} \\
$A$ & Emission-line peak amplitude & \S\ref{sssec:emline_model} \\
$\lambda^*$ & Observed-frame emission-line center & \S\ref{sssec:emline_model} \\
$\Delta v$ & Line-of-sight velocity shift (distinct from the SPS pixel size above) & \S\ref{sssec:emline_model} \\
$\sigma$ & Emission-line width (physical/Doppler), in \kms{}, the emission-line analog of \vdisp{} & \S\ref{sssec:emline_model} \\
$\sigma_\ell$ & Emission-line width in log-$\lambda$ units, $\sigma_\ell=\sigma/c$ & \S\ref{sssec:emline_model} \\
$\Phi$ & Total integrated emission-line flux & \S\ref{sssec:emline_model} \\
$n_\mathrm{line}$ & Number of emission lines (\S\ref{ssec:emlinelist}) within the observed-frame wavelength range at redshift $z$ & \S\ref{sssec:emline_model} \\
$N_\mathrm{dof}$, $\Delta N_\mathrm{dof}$ & Degrees of freedom / change thereof & \S\ref{sssec:kinematic} \\
$N_\sigma$ & Pixel-window half-width, in units of $\sigma_\mathrm{eff}$ & \S\ref{sssec:kinematic} \\
$\sigma_\mathrm{broad}$, $\sigma_\mathrm{narrow}$ & Fitted broad-/narrow-component line width & \S\ref{sssec:kinematic} \\
\cutinhead{Resolution-Matrix Deconvolution}
$c$ & Speed of light & \S\ref{sssec:emline_model} \\
$\mathbf{R}$ & DESI wavelength-dependent resolution matrix & \S\ref{sssec:emline_model} \\
$\mathbf{M}$ & Gaussian design matrix used to deconvolve the resolution matrix & \S\ref{sssec:emline_model} \\
$\mathbf{W}$ & Combined resolution/deconvolution operator, $\mathbf{W}=\mathbf{R}\mathbf{M}^{-1}$ & \S\ref{sssec:emline_model} \\
$\mathbf{h}$ & $\sigma_G$-pre-convolved line profile at pixel centers & \S\ref{sssec:emline_model} \\
$G(\cdot;0,\sigma_G)$ & Gaussian pre-convolution kernel & \S\ref{sssec:emline_model} \\
$\sigma_G$, $\sigma_{G,\ell}$, $\sigma_\mathrm{eff}$ & Pre-convolution kernel width and its log-$\lambda$/effective forms & \S\ref{sssec:emline_model} \\
$x_j$ & Log-wavelength offset from line center at pixel $j$ & \S\ref{sssec:emline_model} \\
$F_j$ & Pre-convolved single-line profile at pixel $j$ & \S\ref{sssec:emline_model} \\
$\mathbf{F}$ & Vector of $F_j$ over all emission lines & \S\ref{sssec:emline_model} \\
$\mathbf{E}$ & Pixelized emission-line model, $\mathbf{E}=\mathbf{W}\mathbf{F}$ & \S\ref{sssec:emline_model} \\
$E_j$ & Emission-line model evaluated at pixel $j$ & \S\ref{sssec:emline_model} \\
$r_i$ & Continuum-subtracted residual spectrum & \S\ref{sssec:emline_model} \\
\cutinhead{Derived Quantities and Uncertainties}
EW & Equivalent width & \S\ref{sssec:lines} \\
$D_n(4000)$ & 4000~\AA{} break index & \S\ref{sssec:lines} \\
$I_A$ & Inverse variance of the flux density at a line center & \S\ref{sssec:lines} \\
$N_\mathrm{monte}$ & Number of Monte Carlo noise realizations & \S\ref{ssec:uncertainties} \\
$\sigma_\theta$ & Monte Carlo uncertainty on a fitted quantity $\theta$ & \S\ref{ssec:uncertainties} \\
\enddata
\tablecomments{Symbols are grouped by role rather than strictly by order of
  appearance. Several symbols carry more than one meaning in
  Section~\ref{sec:modeling}: $C$, $D$, $F$, $M$, and $\mathcal{R}$/$\mathbf{R}$
  each denote two or more distinct quantities, distinguished by typeface
  (roman/bold/calligraphic) and subscript; $\Delta v$ denotes both the SPS
  template pixel size (\S\ref{ssec:templates}) and the emission-line velocity
  shift (\S\ref{sssec:emline_model}), which are unrelated quantities that
  happen to share a symbol.}
\end{deluxetable*}


A key algorithmic choice we make in \fastspec{} is to fit $C$ and $E$
sequentially rather than jointly, subtracting the best-fitting
continuum model from the data before fitting the emission lines
(\S\ref{ssec:emlines}). We adopt this two-stage design primarily for
flexibility: because $C$ and $E$ draw on different spectral
information---the overall continuum shape and broadband photometry for
$C$, (relatively) narrow ``residual'' features for $E$---decoupling
them lets us fit the lines robustly even when the continuum model is a
poor fit to the data due to, for example, residual calibration errors
(\S\ref{sssec:smooth}).

In the remainder of this section we begin by describing the fitting
pipeline's two primary model inputs, the SPS template library
(\S\ref{ssec:templates}) and the emission-line list
(\S\ref{ssec:emlinelist}). We then describe the three sequential
algorithmic stages \fastspec{} executes in turn---adaptive
emission-line masking via fast, preliminary fits in wavelength
patches (\S\ref{ssec:patches}), stellar continuum modeling
(\S\ref{ssec:continuum}), and emission-line modeling
(\S\ref{ssec:emlines}). We close by describing the physical
quantities \fastspec{} derives from these fits
(\S\ref{ssec:quantities}), how we estimate their uncertainties
(\S\ref{ssec:uncertainties}), and the additional fitting modes,
including a ``photometry-only'' mode called \fastphot, built on this
same machinery (\S\ref{sec:modes}).

\subsection{SPS Template Library}
\label{ssec:templates}

\begin{figure*}
\includegraphics[width=6.5in]{figures/sps-models.pdf}
\caption{The five SPS templates spanning \fastspecfit's lookback-time
  age bins (\S\ref{ssec:templates}), computed at solar metallicity
  ($\log_{10}Z/Z_\odot=0$) with $\tauv=0.1$ and normalized to a
  $10^{10}\,\msun$ galaxy at $z=0.1$. Each curve corresponds to one
  age bin, assuming a constant star-formation rate over its duration,
  ranging from 0--30~Myr (youngest) to 11.6--13.7~Gyr (oldest); the
  SED shifts from UV-bright continua with prominent nebular emission
  in the youngest populations to red, featureless spectra in the
  oldest. The gray band marks the DESI optical spectral range
  (0.36--0.98~$\mu$m); the curves along the bottom of the panel show
  the DECam $grz$ and WISE W1-W4 broadband filter throughputs that
  \fastspecfit{} includes in the simultaneous spectrophotometric
  fit. \label{fig:sps}}
\end{figure*}

We generate SPS templates using Flexible Stellar Population
Synthesis\footnote{\url{https://github.com/moustakas/fsps}}
\citep[FSPS;][]{conroy09a,conroy10b} via a fork of the
\texttt{python-fsps}\footnote{\url{https://github.com/moustakas/python-fsps}}
interface. We construct these models using a custom version of the
\texttt{C3K} theoretical spectral library,
\texttt{C3K\_R3K}\footnote{\url{https://github.com/moustakas/c3k}},
the MIST stellar isochrones \citep{choi16a,dotter16a}, and a
\citet{chabrier03a} IMF. The \texttt{C3K\_R3K} library spans 100~\AA{}
to 1000~$\mu$m with $\mathcal{R}_{\mathrm{C3K}} \approx 7065$
($\lambda/\mathrm{FWHM}=3000$) in the rest-frame optical
(3500--9800~\AA) and constant-velocity $\Delta v=50$~\kms{} pixels. At
shorter and longer wavelengths our models have lower resolution and
sparser wavelength sampling, which keeps the UV-to-IR template
lightweight for synthetic photometry and dust-emission modeling
(\S\ref{sssec:fitting}) without wasting resolution where the
spectroscopic fit does not need it.

Following \citet{leja17a}, we build five ``continuity age'' basis
templates spanning fixed lookback-time bins, each assuming a constant
star-formation rate over its duration. The two youngest bins span
0--30~Myr and 30--100~Myr (with fixed widths of
$\Delta\log_{10}(t/\mathrm{yr}) = 7.48$ and $8.0$, respectively),
while the three oldest bins are logarithmically spaced between 100~Myr
and 13.7~Gyr with binning 0.1--1.1~Gyr, 1.1--11.6~Gyr, and
11.6--13.7~Gyr (Figure~\ref{fig:sps}). We generate templates at fixed
solar metallicity, $\log_{10}(Z/Z_\odot) = 0$ (with no
$\alpha$-element enhancement), in two parallel variants---one
including and one excluding nebular emission (continuum and lines)
from the FSPS built-in nebular photoionization grid \citep{byler17a}.

We index the resulting rest-frame templates $S_n(\lambda)$, $n =
1,\ldots,N$ (with $N=5$), by star-formation history (SFH) age bin and
nebular-emission variant, and normalize each to the flux density
produced at 10~pc by one solar mass of stars formed at the tabulated
lookback time,
\begin{equation}
  S_n(\lambda) =
  \frac{L_\odot}{4\pi(10\,\mathrm{pc})^2}\,\frac{L_{\lambda,n}}{L_\odot},
  \label{eq:fnorm}
\end{equation}
in units of
$\mathrm{erg\,s^{-1}\,cm^{-2}\,\text{\AA}^{-1}\,\msun^{-1}}$, so that
each template's fitted coefficient $c_n \geq 0$
(\S\ref{sssec:fitting}) carries units of \msun.

In addition, we generate a single infrared dust emission template
$S_\mathrm{dust}(\lambda)$ from \citet{draine07a}. This model is
parameterized by the polycyclic aromatic hydrocarbon (PAH) mass
fraction, $q_\mathrm{PAH}$, the minimum starlight intensity heating
the diffuse ISM dust component $U_\mathrm{min}$, and the fraction
$\gamma$ of dust mass exposed to a power-law distribution of starlight
intensities (up to a fixed maximum $U_\mathrm{max}$) representing
photodissociation regions. Because most DESI targets do not have
adequately deep IR photometry to constrain a large number of dust
emission parameters, we adopt fiducial values of $q_\mathrm{PAH} =
3.5$, $U_\mathrm{min} = 1.0$, and $\gamma = 0.01$, and normalize our
IR template to unit bolometric luminosity.

There is no AGN-like emission (either continuum or emission lines) in
these templates, a shortcoming which should be kept in mind throughout
this paper. However, \S\ref{ssec:improvements} discusses some of the
planned improvements for \fastspecfit{}, including a QSO/AGN fitting
mode.

%\begin{deluxetable}{ccl}
\tablecaption{SPS Model Parameters\label{table:sps}}
\tablewidth{0pt}
\tablehead{
\colhead{Parameter} & 
\colhead{Values} &
\colhead{Comments} 
}
\startdata
$\log_{10} (Z/Z_{\odot})$ & -1.0, 0.0, 0.3 & Stellar metallicity. \\
$A_{V}$ & 0.0, 0.021668548, & Dust attenuation. \\
\enddata
\tablecomments{All models assume a fixed \citet{chabrier03a} initial mass function from $0.1-100~\mathcal{M}_{\odot}$ and a \citet{charlot00a} dust attenuation curve, $A(\lambda)\propto\lambda^{\alpha}$, with a fixed $\alpha=-0.7$ slope.}
%\tablenotemark{a}{}
\end{deluxetable}


\subsection{Emission-Line List}
\label{ssec:emlinelist}

The default \fastspecfit{} line-list, \texttt{data/emlines.ecsv},
includes 46 emission lines, spanning rest-frame
Ly$\alpha$~$\lambda$1215 through [\ion{S}{3}]~$\lambda$9532 (see
Table~\ref{table:emlines}). We select this list to span the full range
of galaxy and QSO spectral types over the DESI redshift range,
targeting the strongest and most diagnostically useful transitions
rather than every possible detectable line; the resulting set supports
target selection and scientific follow-up across a broad range of
object classes, from star-forming galaxies to high-redshift
quasars.

The adopted vacuum rest-frame wavelengths in Table~\ref{table:emlines}
come from the NIST Atomic Spectra
Database\footnote{\url{https://physics.nist.gov/PhysRefData/ASD/lines_form.html}}
for the Balmer series, the helium lines, the \mgiidoublet{} doublet,
and most forbidden transitions; the UV broad-line wavelengths come
from \citet{vanden-berk01a}; and wavelengths for \niidoublet,
\oiiidoublet, \nii~$\lambda5755$, and \oiii~$\lambda4363$ come from
the CHIANTI atomic
database.\footnote{\url{https://db.chiantidatabase.org}}

We organize these lines into four general categories---UV/QSO lines,
narrow Balmer/helium lines, broad Balmer/helium lines, and narrow
forbidden lines---which are used to set kinematic fitting priors (see
\S\ref{sssec:kinematic}). Table~\ref{table:emlines} also lists a set
of 23 spectrally adjacent ``patches'', and a Boolean bit (``strong'')
which flags which lines, if present, are most likely to be
well-detected in a given spectrum. We describe these latter quantities
in detail in the next subsection (\S\ref{ssec:patches}).

\todo{The table should add a comment about the ``UV/QSO''
  categorization and also explain the \texttt{isbroad} designation not
  shown here (but implicitly used in the next section).}

\startlongtable
\begin{deluxetable}{llccc}
\tablecaption{Emission Lines\label{table:emlines}}
\tablewidth{0pt}
\tablehead{
\colhead{Emission} &
\colhead{Line} &
\colhead{Rest} &
\colhead{Initial} &
\colhead{} \\
\colhead{Line} &
\colhead{Prefix\tablenotemark{a}} &
\colhead{Wavelength\tablenotemark{b}} &
\colhead{Patch\tablenotemark{c}} &
\colhead{``Strong''\tablenotemark{d}}
}
\startdata
\cutinhead{UV/QSO Lines}
Ly$\alpha$                        & \texttt{LYALPHA}        & 1215.670   & a & \checkmark \\
{[\ion{N}{5}]}~$\lambda$1240      & \texttt{NV\_1240}       & 1240.14    & a & \checkmark \\
{\ion{O}{1}}~$\lambda$1304        & \texttt{OI\_1304}       & 1304.35    & b &            \\
{\ion{Si}{4}}~$\lambda$1396       & \texttt{SILIV\_1396}    & 1396.76    & c &            \\
{\ion{C}{4}}~$\lambda$1549        & \texttt{CIV\_1549}      & 1549.06    & d & \checkmark \\
{\ion{He}{2}}~$\lambda$1640       & \texttt{HEII\_1640}     & 1640.42    & e &            \\
{\ion{Al}{3}}~$\lambda$1857       & \texttt{ALIII\_1857}    & 1857.40    & f &            \\
{\ion{Si}{3}}]~$\lambda$1892      & \texttt{SILIII\_1892}   & 1892.030   & f & \checkmark \\
{\ion{C}{3}}]~$\lambda$1908       & \texttt{CIII\_1908}     & 1908.734   & f & \checkmark \\
{\ion{Mg}{2}}~$\lambda$2796       & \texttt{MGII\_2796}     & 2796.3511  & g & \checkmark \\
{\ion{Mg}{2}}~$\lambda$2803       & \texttt{MGII\_2803}     & 2803.5324  & g & \checkmark \\
{\ion{He}{2}}~$\lambda$4686       & \texttt{HEII\_4686}     & 4687.02    & o &            \\
\cutinhead{Narrow Balmer/Helium Lines}
H6                                & \texttt{H6}             & 3890.1576  & j &            \\
H$\epsilon$                       & \texttt{HEPSILON}       & 3971.195   & k &            \\
H$\delta$                         & \texttt{HDELTA}         & 4102.892   & l &            \\
H$\gamma$                         & \texttt{HGAMMA}         & 4341.684   & m & \checkmark \\
{\ion{He}{1}}~$\lambda$4471       & \texttt{HEI\_4471}      & 4472.7290  & n &            \\
H$\beta$                          & \texttt{HBETA}          & 4862.683   & p & \checkmark \\
{\ion{He}{1}}~$\lambda$5876       & \texttt{HEI\_5876}      & 5877.2690  & r &            \\
H$\alpha$                         & \texttt{HALPHA}         & 6564.613   & t & \checkmark \\
\cutinhead{Broad Balmer Lines}
H$6_{\rm b}$                      & \texttt{H6\_BROAD}       & 3890.1576  & j &            \\
H$\epsilon_{\rm b}$               & \texttt{HEPSILON\_BROAD} & 3971.195   & k &            \\
H$\delta_{\rm b}$                 & \texttt{HDELTA\_BROAD}   & 4102.892   & l &            \\
H$\gamma_{\rm b}$                 & \texttt{HGAMMA\_BROAD}   & 4341.684   & m & \checkmark \\
H$\beta_{\rm b}$                  & \texttt{HBETA\_BROAD}    & 4862.683   & p & \checkmark \\
H$\alpha_{\rm b}$                 & \texttt{HALPHA\_BROAD}   & 6564.613   & t & \checkmark \\
\cutinhead{Narrow Forbidden Lines}
{[\ion{Ne}{5}]}~$\lambda$3346     & \texttt{NEV\_3346}      & 3346.7828  & h &            \\
{[\ion{Ne}{5}]}~$\lambda$3426     & \texttt{NEV\_3426}      & 3426.8637  & h &            \\
{[\ion{O}{2}]}~$\lambda$3726      & \texttt{OII\_3726}      & 3727.0919  & i & \checkmark \\
{[\ion{O}{2}]}~$\lambda$3729      & \texttt{OII\_3729}      & 3729.8750  & i & \checkmark \\
{[\ion{Ne}{3}]}~$\lambda$3869     & \texttt{NEIII\_3869}    & 3869.8611  & j &            \\
{[\ion{O}{3}]}~$\lambda$4363      & \texttt{OIII\_4363}     & 4364.4351  & m &            \\
{[\ion{O}{3}]}~$\lambda$4959      & \texttt{OIII\_4959}     & 4960.2937  & p & \checkmark \\
{[\ion{O}{3}]}~$\lambda$5007      & \texttt{OIII\_5007}     & 5008.2383  & p & \checkmark \\
{[\ion{N}{2}]}~$\lambda$5755      & \texttt{NII\_5755}      & 5756.1887  & q &            \\
{[\ion{O}{1}]}~$\lambda$6300      & \texttt{OI\_6300}       & 6302.0435  & s &            \\
{[\ion{S}{3}]}~$\lambda$6312      & \texttt{SIII\_6312}     & 6313.8062  & s &            \\
{[\ion{N}{2}]}~$\lambda$6548      & \texttt{NII\_6548}      & 6549.8578  & t & \checkmark \\
{[\ion{N}{2}]}~$\lambda$6584      & \texttt{NII\_6584}      & 6585.2696  & t & \checkmark \\
{[\ion{S}{2}]}~$\lambda$6716      & \texttt{SII\_6716}      & 6718.2913  & u & \checkmark \\
{[\ion{S}{2}]}~$\lambda$6731      & \texttt{SII\_6731}      & 6732.6705  & u & \checkmark \\
{[\ion{Ar}{3}]}~$\lambda$7135     & \texttt{ARIII\_7135}    & 7137.77    & v &            \\
{[\ion{O}{2}]}~$\lambda$7320      & \texttt{OII\_7320}      & 7321.731   & v &            \\
{[\ion{O}{2}]}~$\lambda$7330      & \texttt{OII\_7330}      & 7332.199   & v &            \\
{[\ion{S}{3}]}~$\lambda$9069      & \texttt{SIII\_9069}     & 9071.103   & w &            \\
{[\ion{S}{3}]}~$\lambda$9532      & \texttt{SIII\_9532}     & 9533.227   & w &            \\
\enddata

\tablenotemark{a}{Data-model prefix for every output quantity we
  measure for this line (\S\ref{sssec:lines}), e.g.,
  \texttt{OIII\_5007\_FLUX}}
\tablenotemark{b}{Vacuum wavelengths in \AA{} come from the NIST Atomic Spectra
Database, the CHIANTI atomic database, or from \citet{vanden-berk01a};
see \S\ref{ssec:emlinelist} for more details.}
\tablenotemark{c}{Initial assignment of spectrally adjacent lines which are jointly fit by the line-masker (see \S\ref{ssec:patches}).}
\tablenotemark{d}{Lines fit in the line-masker's initial patch-fitting pass (\S\ref{ssec:patches}) to estimate robust kinematics; if present, these are the lines most likely to be well detected.}
\end{deluxetable}


%Twelve UV and AGN broad lines---from Ly$\alpha$ through
%\ion{He}{2}~$\lambda$4686---carry broad, independently fitted
%kinematics appropriate for quasar and AGN broad-line emission. Eight
%Balmer and helium lines (H6 through H$\alpha$, together with
%\ion{He}{1}~$\lambda\lambda$4471, 5876) receive narrow kinematic
%components, and the six Balmer lines (H6 through H$\alpha$)
%additionally receive independent broad components, enabling
%\fastspecfit{} to separate star-forming and AGN contributions to the
%Balmer emission. The remaining 20 lines are forbidden transitions,
%each fitted with a single narrow Gaussian. We describe the kinematic
%groupings and amplitude constraints that link these lines during
%fitting in \S\ref{sssec:kinematic}.

\subsection{Adaptive Emission-Line Masking}
\label{ssec:patches}

In order to fit the stellar continuum robustly, we must first identify
pixels which could be affected by emission lines. The challenge is
that DESI targets span an enormous range of emission-line properties,
from purely absorption-line quiescent galaxies with no measurable
line-emission at all to extreme starbursts and broad-line AGN whose
emission dominates large stretches of the observed spectrum. A single,
fixed masking window sized for the strongest expected lines would
over-mask quiescent spectra and discard usable continuum information,
while a window sized for narrow lines found in typical star-forming
galaxies would under-mask galaxies and QSOs with broad emission lines
and bias the stellar continuum fit.

Therefore, before fitting the stellar continuum
(\S\ref{sssec:fitting}), we run a lightweight, per-object fitting
pass---the \textit{line masker}---that determines whether a spectrum
contains statistically significant line-emission, and, if so, the
approximate velocity widths and velocity shifts of those
lines. Crucially, the line-masker also performs an initial
determination of whether the spectrum contains broad Balmer
line-emission (although the statistically defensible decision is made
in the final fitting pass, described in \S\ref{sssec:kinematic}).

\begin{figure*}
\includegraphics[width=6.5in]{figures/linemasker.pdf}
\caption{Illustration of the line-masking algorithm for an example BGS
  target at $z=0.3292$ (\texttt{TARGETID=39627670102744387} from the
  main survey bright-time program, appearing on HEALPix
  17366). \todo{Consider adding an image cutout.} \textit{Top:} the
  full per-camera DESI spectrum (B, R, Z in blue, green, and red,
  respectively) together with the per-patch best-fitting local linear
  continuum and Gaussian model line-profiles (darker blue, green, and
  red in B, R, and Z, respectively) for all the candidate strong lines
  in the observed spectral range. \textit{Bottom:} Zoomed-in panels
  which more clearly show the per-patch fitting results. In each
  panel, camera-colored pixels mark the pixels in the final line-mask
  (and therefore excluded from the subsequent continuum-fit), while
  light-gray pixels are the surrounding continuum used for that
  patch's local linear fit. The darker curves show the combined
  best-fit line-plus-continuum model, while the dashed black line and
  surrounding dotted band show the local linear continuum and its
  $\pm1\sigma$ noise envelope. Finally, the legends indicate the S/N
  value inferred for each line and the patch label. For this
  particular object, a significant broad \ha{} Balmer line was
  correctly identified, resulting in patches ``m'' and ``p'' and
  patches ``t'' and ``u'' merging into larger patches ``mp'' and
  ``tu'', respectively, and more aggressive (conservative)
  line-masking. \label{fig:linemasker}}
\end{figure*}

%Each patch groups lines whose wavelength windows overlap or nearly
%touch; for example, one patch contains \hb{} together with
%\oiiidoublet, and another groups \ha{} with \niidoublet.

%Because the members of each kinematic groups span many patches, all
%the this is a single joint fit of line amplitudes and per-patch linear
%continua across all active patches, not a set of independent per-patch
%fits.

First, we activate all the spectral patches (see
Table~\ref{table:emlines}) with at least one line in the observed
spectral range and fit all the ``strong'' lines with the narrow-only
kinematic model, together with a simple linear local continuum in each
patch. The lines we initially fit are simply the subset of lines
which, if present, are most likely to yield useful kinematic
constraints. In the narrow-only kinematic model, all the forbidden
lines are tied together into a single kinematic group (shared velocity
width and velocity center) anchored on \oiiilam, while all the narrow
Balmer and \ion{He}{1} lines are anchored on \ha{} in a second
kinematic group (see \S\ref{sssec:kinematic} for more details). Then,
we define a pixel window of $\pm2.5\,\sigma_\mathrm{line}$ centered on
each line's expected observed-frame wavelength, where
$\sigma_\mathrm{line}$ is that line's current velocity-width
estimate. We iteratively update $\sigma_\mathrm{line}$, beginning with
150~\kms{} and 3000~\kms{} for narrow and (candidate) broad lines,
respectively, but always floored at a minimum line-width of 50~\kms,
and grow a surrounding continuum window until at least 11 unmasked
pixels are available. When this adaptive growth causes two spectrally
adjacent patches' continuum windows to overlap by more than 11 pixels,
we merge them into a single patch, sharing one local continuum and the
union of both patches' lines.

%We run two such iterations in total, updating the line-pixel windows
%with the best-fit velocity shifts and widths before the second pass.

%for example, the merged \ha+\nii+\sii{} complex shown
%in Figure~\ref{fig:linemasker}.

%Table~\ref{table:emlines} lists the set of 46 emission lines in
%\fastspecfit's default line-list, the set of 23 spectrally adjacent
%``patches'' each line belongs to, and a Boolean bit (``strong'') which
%flags which lines, if present, are most likely to be well-detected in
%a given spectrum. We emphasize that the differentiation between
%``strong'' and ``weak'' lines at this point in the pipeline is
%qualitative; the strong lines, if present, are simply the subset of
%lines which are most likely to yield useful kinematic constraints
%(line-widths and velocity shifts), while the subset of well-detected
%lines in the full line-list are used in the final mask-building step,
%as we describe below.

Next, if any broad Balmer line falls within the observed spectral
range, we repeat the procedure with the narrow+broad model,
initializing all the broad Balmer line-widths at 1000~\kms. In this
fitting mode, all the ``strong'' broad Balmer lines are kinematically
tied together, separately from the two previously described
narrow-line kinematic groups (see also \S\ref{sssec:kinematic} for
more details). After this fitting pass, if any broad Balmer line
reaches $\mathrm{S/N} > 1.5$ we adopt the narrow+broad line-widths for
masking, and otherwise retain the narrow-only widths.

With the robustly estimated line-widths, we can now build the complete
line-mask. Iterating over all the lines in the observed spectral range
(not just the initial set of candidate strong lines), we robustly
estimate each line's signal-to-noise ratio by comparing its peak
amplitude against the variance of the pixels in its surrounding local
continuum, and discard ``weak'' lines with $\mathrm{S/N} < 3.5$ from
further consideration, in order to avoid over-masking the continuum
spectrum. For the surviving lines, we adopt a
$\pm5\,\sigma_\mathrm{line}$ window, centered on each line's expected
wavelength and sized using its converged velocity width and velocity
shift, to build the final pixel mask.

%We keep this pass deliberately fast and approximate---fitting a simple
%linear local continuum per patch rather than the full SPS template,
%and adopting loose convergence tolerances---since its only job is to
%hand the main continuum and emission-line fits that follow
%(\S\ref{ssec:continuum}, \S\ref{ssec:emlines}) accurate-enough masks
%and kinematic initial guesses.

Figure~\ref{fig:linemasker} illustrates this process for a single BGS
target from the bright program of the main survey,
\texttt{TARGETID}=39627670102744387, at $z=0.3292$. Two pairs of
spectrally adjacent patches are merged in this example: the ``m''
(H$\gamma$) and ``p'' (\hb{} and \oiiidoublet) patches merge into a
single ``mp'' patch, and the ``t'' (\ha{} and \niidoublet) and ``u''
(\siidoublet) patches merge into ``tu'', because their adaptive
continuum windows overlap by more than 11 pixels. The line-masker also
detects broad Balmer line-emission in this target, so the
corresponding patches are masked more conservatively, using the
narrow+broad kinematic model described above.

%This object-by-object sensitivity to broad-line emission is precisely
%what enables the robust identification of very low-mass black holes in
%dwarf galaxies (e.g., \citealt{pucha25a}).

\subsection{Stellar Continuum Model}
\label{ssec:continuum}

This section describes the steps \fastspec{} uses to find the
best-fitting continuum model, $C(\lambda;\theta_C)$. Because speed and
flexibility are paramount, the default fitting mode uses a minimal
parameter set, $\theta_C = \{z,\mathbf{c},\tauv,\vdisp\}$, where $z$
is the redshift, $\mathbf{c}$ are the non-negative template
coefficients, $\tauv$ is the $V$-band dust optical depth, and $\vdisp$
is the stellar velocity dispersion. We emphasize, however, that the
fitting infrastructure can be easily extended to include additional
parameters (e.g., a two-component dust model or QSO model spectra),
which we discuss, in part, in \S~\ref{ssec:improvements}.

For each object, we determine $\theta_C$ in three stages. First, we
measure $\vdisp$ from stellar absorption features
(\S\ref{sssec:vdisp}); then we derive a scalar aperture correction
that places the DESI fiber spectrum on the total-flux scale of the
broadband photometry (\S\ref{sssec:apercorr}); and finally we jointly
fit for dust attenuation $\tauv$ and the template coefficients
$\mathbf{c}$, holding $z$ fixed at its input (Redrock) value
(\S\ref{sssec:fitting}). An optional nonparametric smooth-continuum
correction then removes residual low-level systematics
(\S\ref{sssec:smooth}) before the spectrum passes to the emission-line
fitter.

\subsubsection{Stellar Velocity Dispersion}
\label{sssec:vdisp}

The line-of-sight stellar velocity dispersion \vdisp{} traces the
depth of a galaxy's gravitational potential well; it is a key
observable for several galaxy scaling relations---the Fundamental
Plane, the black hole mass--\vdisp{} relation, and the
halo-mass--stellar-mass relation---and a fundamental ingredient for
most theoretical models of galaxy formation \citep[and references
  therein]{renzini06a, blanton09a, kormendy13a, vegetti24a}. A
reliable measurement, however, requires adequate rest-frame wavelength
coverage and relatively high signal-to-noise spectroscopy of the
absorption-line features most sensitive to dynamical broadening (e.g.,
Ca~H\&K, the $G$-band, H$\beta$, Mg~$b$). In DESI, these conditions
are generally met only for BGS and LRG targets; however, because
\fastspec{} does not use targeting information as a prior, we attempt
to measure \vdisp{} for every spectroscopic target.

We determine whether a direct measurement of \vdisp{} is possible from
two conditions: rest-frame wavelength coverage and ``fit
significance.'' First, we require a given target to have at least
1100~\AA{} of (unmasked) rest-frame wavelength coverage between 3800
and 4900~\AA, bridging Ca~H\&K, the 4000-\AA{} break, and H$\beta$. If
this condition is satisfied, then we perform the actual fit using all
unmasked pixels between 3800--6000~\AA{} (in the rest-frame). In
practice, this requirement means that we can only estimate \vdisp{}
for galaxies up to $z\approx1$. Second, we determine fit significance
by performing a fast $\chi^2$ grid scan of six trial
velocity-broadening widths spanning $[50,\,500]$~\kms{},
simultaneously solving for \tauv{} and the template coefficients by
linear least-squares at each grid point (see
\S\ref{sssec:fitting}). If this scan yields a peak-to-peak
$\Delta\chi^2 \geq 25$ (and with the minimum value away from either
boundary) then we proceed with a direct measurement. In general,
$\Delta\chi^2 < 25$ implies that the spectrum has inadequate
signal-to-noise ratio in the stellar continuum (e.g., a faint ELG
target) to accommodate \vdisp{} as an additional free parameter.

Most DESI targets in DR2 fail one or both of the preceding conditions;
however, we would still like to use an approximately correct value of
\vdisp. In this case, we estimate \vdisp{} from the
stellar-mass--velocity-dispersion relation,
\begin{equation}
  \log_{10}(\vdisp/\kms) = a + b\,(\logmstar - 11),
  \label{eq:sigma_mstar}
\end{equation}
with $(a,b) = (2.30,\,0.25)$ (\todo{need a ref}) using a fast,
preliminary estimate of the stellar mass, \Mstar{} (see
\S\ref{sssec:fitting}). If our estimate of \Mstar{} fails---for
whatever reason---we instead adopt a fallback $\vdisp = 150$~\kms{}
(\texttt{VDISP\_NOMINAL}), close to the characteristic ``knee'' of the
velocity dispersion function \citep{bezanson11a, taylor22a}. In either
fallback case we set \texttt{VDISP\_IVAR}~$=0$ to indicate that
\vdisp{} has not been independently measured.

When a direct measurement is warranted, we determine \vdisp{} by
finding the optimal Gaussian kernel dispersion,
$\vdisp^\mathrm{kern} \in [50,\,500]$~\kms{}, using
\texttt{scipy.optimize.least\_squares} with a Trust-Region Reflective
minimization algorithm \citep{branch99a}, fitting for
$\vdisp^\mathrm{kern}$ and \tauv{} simultaneously over the same range
as above. We then recover the physical velocity dispersion by adding
back the intrinsic resolution of our SPS models in quadrature,
\begin{equation}
  \vdisp = \sqrt{(\vdisp^\mathrm{kern})^2 + \sigma_\mathrm{C3K}^2},
  \label{eq:vdisptot}
\end{equation}
where $\sigma_\mathrm{C3K} = c/\mathcal{R}_{\mathrm{C3K}}=42$~\kms{}
(\S\ref{ssec:templates}).

\subsubsection{Aperture Correction}
\label{sssec:apercorr}

By default, \fastspec{} fits the spectrum ($0.36-0.98$\,\micron) of
each extragalactic DESI target simultaneously with its corresponding
seven-band ($grz$W1--W4) Legacy Surveys integrated photometry
($0.47-22$\,\micron; \S\ref{ssec:photometry}). Fitting all the
available spectrophotometry over the broadest possible wavelength
baseline helps constrain physical properties like dust attenuation,
and ties the derived stellar masses, star-formation rates, and
rest-frame photometry to the total light of each source rather than to
the light captured by the $1\farcs5$-diameter DESI fiber. However, we
have to account for the fact that the spectrum and photometry are
measured over different spatial scales before combining them.

So-called aperture effects (or aperture bias) are extremely difficult
to treat in detail because galaxies exhibit genuine radial gradients
(or, more generally, spatial inhomogeneities) in color, stellar
populations, age, metallicity, and more \citep[and references
  therein]{sanchez20a}. In addition, the amount of light the DESI
fiber captures depends on each galaxy's apparent (angular) size on the
sky, and therefore on redshift through the angular diameter distance
\citep{hogg99a}. These effects conspire to make aperture bias a strong
function of target class, affecting low-redshift BGS targets the most
and higher-redshift, angularly compact, and more
spectrophotometrically homogeneous LRG and ELG targets the least.

Nevertheless, in the current version of \fastspec{} we use the
simplest possible model to account for aperture bias: we assume that
the spectral shape within the fiber is representative of the galaxy as
a whole, independent of wavelength or the magnitude of the aperture
correction. We derive this correction by synthesizing $grz$ photometry
from a fast initial fit to the DESI spectrum to get $C_k^\mathrm{fib}$
for $k \in \{g,\,r,\,z\}$, and computing the per-band multiplicative
factor
\begin{equation}
  \eta_k = \frac{D_k}{C_k^\mathrm{fib}},
  \label{eq:apercorr}
\end{equation}
where $D_k$ is the total Tractor flux of the object in bandpass
$k$. Finally, we compute the median scalar aperture correction,
$\bar{\eta}=\mathrm{median}(\eta_g,\,\eta_r,\,\eta_z)$, and apply the
resulting multiplicative factor (falling back to $\bar{\eta}=1.0$ if
any part of the preceding analysis fails) during the full
spectrophotometric fit (\S\ref{sssec:fitting}). \todo{Do we need to
  say any more? Or point to a discussion of the limitations of this
  correction?}

%large band-to-band scatter indicates strong color gradients within the
%fiber aperture. We treat $\bar{\eta}$ as exact and do not propagate
%additional uncertainty from the aperture-correction step into the fit
%or the output parameter uncertainties (\S\ref{ssec:uncertainties}).
%which is generally $\gtrsim 1$ but can scatter slightly below unity.
%We adopt the scalar correction $\bar{\eta} =
%\mathrm{median}(\eta_g,\,\eta_r,\,\eta_z)$, which by construction is
%wavelength-independent, and scale the fiber spectrum and its
%inverse-variance array by $\bar{\eta}$ before the joint
%spectrophotometric fit, placing the spectrum on the total-flux scale;
%output model spectra are then divided by $\bar{\eta}$ to recover
%fiber-aperture units.
%
%In \S\ref{sec:modes}, we illustrate the effect of disabling this
%aperture correction.

\subsubsection{Fitting the Stellar Continuum}
\label{sssec:fitting}

To fit the stellar continuum we use a fully forward-modeling approach.
We carry each rest-frame SPS template (\S\ref{ssec:templates}) through
a sequence of physical and instrumental transformations---dust
attenuation, redshifting, and instrumental broadening (or filter
convolution)---into a prediction that can be compared directly to the
observed spectrophotometry, and then solve for the template weights
and dust attenuation (at fixed \vdisp; see \S\ref{sssec:vdisp}) that
minimize the inverse-variance-weighted residuals.

First, we model the attenuation of each template as a
uniform foreground screen,
\begin{equation}
  S_n^\mathrm{att}(\lambda) = S_n(\lambda)\,e^{-\tauv k(\lambda)},
  \label{eq:atten}
\end{equation}

\noindent where $S_n(\lambda)$ is the unattenuated template given by
equation~(\ref{eq:fnorm}), $\tauv$ is the $V$-band optical depth, and

\begin{equation}
  k(\lambda) =
  \left(\frac{\lambda}{5500\,\text{\AA}}\right)^{n}
  \label{eq:kl}
\end{equation}

\noindent is the attenuation curve, which we model as a simple power
law with fixed power-law index $n=-0.7$ based on the
\citet{charlot00a} birth-cloud model.  Absorbed photons are re-emitted
in the IR following energy balance: we scale the \citet{draine07a}
dust emission template, $S_\mathrm{dust}(\lambda)$, by the luminosity
absorbed by each template,
\begin{equation}
  L_{\mathrm{abs},n} = \int_0^\infty S_n(\lambda) (1 - e^{-\tauv
    k(\lambda)})\,d\lambda,
  \label{eq:labs}
\end{equation}
so that the total (attenuated plus reradiated) continuum template
conserves bolometric luminosity,
\begin{equation}
  S_n^\mathrm{tot}(\lambda) = S_n^\mathrm{att}(\lambda) +
  L_{\mathrm{abs},n}\,S_\mathrm{dust}(\lambda).
  \label{eq:ftot}
\end{equation}

For an object at redshift $z$ with luminosity distance $D_L$ (in Mpc),
we relate each rest-frame template to its observed-frame counterpart
via
\begin{equation}
  S_n^\mathrm{obs}(\lambda)
  =
  T_\mathrm{IGM}(\lambda,z)\,
  \frac{\mathcal{F}_0\,\mathcal{M}_0}
       {4\pi\,(3.086\times10^{24}\,D_L)^2}
  \,\frac{S_n^\mathrm{tot}(\lambda/(1+z))}{1+z},
  \label{eq:zfactor}
\end{equation}
where $T_\mathrm{IGM}(\lambda,z)$ is the \citet{inoue14a}
intergalactic medium (IGM) transmission function which accounts for
Ly$\alpha$ forest and Lyman-limit attenuation; $\mathcal{F}_0 =
10^{17}$ is a dimensionless scale factor converting the cgs flux
densities of \S\ref{ssec:templates} to our internal flux units of
$10^{-17}$~erg~s$^{-1}$~cm$^{-2}$~\AA$^{-1}$ \citep{guy23a}; and
$\mathcal{M}_0 = 10^{10}\,\msun$ is a fixed mass normalization that
rescales the fitted template coefficients to order unity for a
well-conditioned non-negative least-squares (NNLS) solve (see below).

%We combine the scalar prefactor and IGM transmission into a single
%array (\texttt{zfactors}) evaluated once per object at initialization
%and reused for all template evaluations during fitting.

We restrict the dust optical depth to the physically motivated range
$\tauv \in [0,\,2]$, and by default we include coefficients for all
five SPS age bins (\S\ref{ssec:templates}), even for stellar
populations nominally older than the age of the universe at the
redshift of a given object. Because our age sampling is coarse, the
boundary between physical and unphysical ages is poorly defined, and
excluding bins outright would introduce more redshift-dependent bias
than it removes \citep{leja19a}.

At this point, we have all the ingredients necessary to fit the
stellar continuum. Let $d$, $I$, and $m$ be the observed DESI
spectrum, inverse-variance spectrum, and binary emission-line mask
from \S\ref{ssec:patches}, respectively, each an
$n_\mathrm{pix}$-element array ($m_i=1$ for pixels retained in the fit
and $0$ for masked pixels); and let $D$ and $I_D$ be the
$n_\mathrm{band}=7$ ($grz$W1--W4) photometric fluxes and inverse
variances, respectively. We construct the continuum model
$C(\lambda;\theta_C)$ as the non-negative linear combination of the
observed-frame templates $S_n^\mathrm{obs}(\lambda)$ from
equation~(\ref{eq:zfactor}),
\begin{equation}
  C(\lambda;\theta_C) = \sum_n c_n\,S_n^\mathrm{obs}(\lambda),
  \label{eq:ccont}
\end{equation}
holding \vdisp{} fixed at the value determined in \S\ref{sssec:vdisp}
and solving only for $\tauv$ (which enters through $S_n^\mathrm{obs}$)
and the coefficients $\mathbf{c}$.

To compare this model against the spectroscopic and photometric data,
we transform it in two ways: (1) we broaden each template by \vdisp{}
(\S\ref{sssec:vdisp}), convolve it with the wavelength-dependent
resolution matrix $\mathbf{R}$ (\S\ref{ssec:spectra}), and resample
(using a flux-conserving trapezoidal-resampling algorithm) onto the
same $n_\mathrm{pix}$-pixel wavelength grid as $d$, giving $C_i \equiv
C(\lambda_i;\theta_C)$; and (2) convolve each template with the
$n_\mathrm{band}$ filter response functions (\S\ref{ssec:photometry})
to derive the synthetic photometry $C_k$. Collecting these fully
transformed templates as the columns of a design matrix
$\boldsymbol{\Psi}$, the continuum model at every pixel and bandpass
becomes
\begin{equation}
  \mathbf{C} = \boldsymbol{\Psi}\,\mathbf{c},
  \label{eq:designmatrix}
\end{equation}
\noindent and we determine $\tauv$ and $\mathbf{c}$ by minimizing
\begin{equation}
  \chi^2(\tauv,\mathbf{c})
  =
  \sum_{i=1}^{n_\mathrm{pix}} m_i I_i\Bigl(d_i - \frac{C_i}{\bar\eta}\Bigr)^2
  \;+\;
  \sum_{k=1}^{n_\mathrm{band}} I_{D,k}\bigl(D_k - C_k\bigr)^2,
  \label{eq:chi2cont}
\end{equation}
where $\bar\eta$ is the scalar aperture correction described in
\S\ref{sssec:apercorr}.

Because equation~(\ref{eq:designmatrix}) is linear in the multiple
template coefficients $\mathbf{c}$ but nonlinear only in the single
scalar $\tauv$, rather than solve for both simultaneously we exploit
the separability via variable projection (VARPRO; \citealt{golub73a}).
With this approach, we perform an outer Brent scalar minimization over
$\tauv$, with an inner NNLS solve for $\mathbf{c}$ at each trial
$\tauv$,
\begin{equation}
  \mathbf{c}^* = \arg\min_{\mathbf{c}\,\geq\,0}
  \bigl\|\boldsymbol{\Psi}\,\mathbf{c} - \mathbf{b}\bigr\|^2,
  \label{eq:nnls}
\end{equation}
where $\mathbf{b}$ is the corresponding weighted data vector, and the
mask, inverse variances, and aperture correction enter as row weights
on $\boldsymbol{\Psi}$ and $\mathbf{b}$ alike---$\sqrt{m_i I_i}/\bar\eta$
for spectral rows and $\sqrt{I_{D,k}}$ for photometric rows. Minimizing
$\|\boldsymbol{\Psi}\mathbf{c}-\mathbf{b}\|^2$ is therefore equivalent
to minimizing equation~(\ref{eq:chi2cont}) at fixed $\tauv$. Because
the inner NNLS solve is exact at every trial $\tauv$, VARPRO converges
an order-of-magnitude faster than direct optimization over the full
parameter vector, which would otherwise require finite-differencing
across all template coefficients.

With $\tauv$ and $\mathbf{c}$ now in hand, and $z$ and $\vdisp$ fixed
from the earlier stages, every element of $\theta_C$ is determined,
giving the best-fitting stellar continuum model
$C(\lambda;\theta_C)$. From this model we derive several physical
properties of interest---stellar mass, star-formation rate, and
stellar age---as described in \S\ref{ssec:quantities}.

\subsubsection{Smooth Continuum Correction}
\label{sssec:smooth}

\todo{Does this section need a figure to illustrate the technique?}

Despite its flexibility, our parametric continuum model $C(\lambda)$
(\S\ref{sssec:fitting}) can leave small- or even large-scale
systematic residuals in the continuum-subtracted spectrum. Depending
on their wavelength-dependence and amplitude, leaving these residuals
unaddressed can severely bias the emission-line modeling
(\S\ref{sssec:emline_model}), especially for faint lines.

The unmodeled residuals can be due to instrumental effects, template
mismatch, or both. For example, the noisy boundary between the three
DESI cameras (\S\ref{ssec:spectra}) is especially susceptible to
residual additive (bias-, dark, or sky-subtraction) errors, or to
multiplicative (flux-calibration) errors. Template mismatch,
meanwhile, is most severe for QSOs and broad-line AGN
(\S\ref{ssec:example_qso}). Their observed continuum typically blends
genuine host-galaxy starlight with \ion{Fe}{2} emission in the
rest-frame UV (roughly 1100--3000\,\AA; \citealt{vestergaard01a}) and
a featureless, non-thermal power-law continuum from accretion onto the
supermassive black hole \citep{vanden-berk01a}. Our SPS templates can
represent the host-galaxy light but not the other two, nor cleanly
separate them.

Keeping these effects in mind, the current version of \fastspec{}
treats all such residuals, irrespective of their origin, as a single,
low-level additive correction to $C(\lambda)$. We estimate and remove
these residuals with a nonparametric smooth continuum after
subtracting the best-fitting, line-free continuum model. Working per
camera, we slide a 75-pixel window in 125-pixel steps across the
residual spectrum, applying a $2\sigma$ iterative clip in each window
and computing the robust median and standard deviation. We discard
windows with fewer than 15 unmasked pixels. We then fit a weighted
quadratic spline (\texttt{scipy.interpolate.UnivariateSpline}, $k=2$;
\citealt{virtanen20a}) to the window medians and add the resulting
smooth baseline back to $C(\lambda)$ before emission-line fitting. For
each camera we also record \texttt{SMOOTHCORR\_B},
\texttt{SMOOTHCORR\_R}, and \texttt{SMOOTHCORR\_Z}, the median of the
smooth correction (using all unmasked pixels) divided by the observed
flux, expressed as a percentage---a simple, per-object diagnostic of
the correction's overall amplitude relative to the source itself.

\subsection{Emission-Line Model}
\label{ssec:emlines}

After subtracting the best-fitting continuum model
$C(\lambda;\theta_C)$ from the data $d_i$ (\S\ref{sssec:fitting}),
including the smooth continuum correction (\S\ref{sssec:smooth}), we
are left with a residual spectrum $r_i \equiv d_i - C_i$ (in
\fluxdensity) containing one or more emission lines
(\S\ref{ssec:emlinelist}). In the following three subsections we
describe how we construct and fit the emission-line model,
$E(\lambda;\theta_E)$ (\S\ref{sssec:emline_model}). The full
emission-line parameter set $\theta_E$ consists of three
parameters---an amplitude, velocity offset, and velocity width---for
each of the $n_\mathrm{line}$ lines in our list. We then describe how
we use physically motivated ``kinematic groups'' among different sets
of lines to reduce the number of free parameters in $\theta_E$,
including (optionally) decomposing the Balmer lines into broad and
narrow components (\S\ref{sssec:kinematic}), and conclude with a
summary of all the quantities we measure for each detected (or
undetected) emission line (\S\ref{sssec:lines}).

\subsubsection{Forward Model}
\label{sssec:emline_model}

We forward-model each object's emission-line spectrum as the
non-negative sum of the $n_\mathrm{line}$ emission lines from our line
list (\S\ref{ssec:emlinelist}) that fall within the observed-frame
wavelength range at the Redrock redshift $z$,
\begin{equation}
  E(\lambda;\theta_E) = \sum_{p=1}^{n_\mathrm{line}}
    \phi(\lambda;\,A_p,\lambda_p^*,\sigma_{\ell,p}),
  \label{eq:emline_sum}
\end{equation}
where we model each individual line as a log-Gaussian in observed
wavelength,
\begin{equation}
  \phi(\lambda;\,A,\lambda^*,\sigma_\ell)
  =
  A\exp\!\left(-\frac{(\ln\lambda - \ln\lambda^*)^{2}}{2\sigma_\ell^{2}}\right).
  \label{eq:intrinsic}
\end{equation}
Here, $A$ is the peak amplitude (in \fluxdensity), $\lambda^*$ is the
observed line center in \AA, including both cosmological ($z$) and
peculiar ($\Delta v$) motion of the line-emitting gas,
\begin{equation}
  \lambda^* = \lambda_\mathrm{rest}
    \left(1 + z + \frac{\Delta v}{c}\right),
  \label{eq:lstar}
\end{equation}
and $\sigma_\ell = \sigma/c$ is the line width in log-$\lambda$ units,
where $\sigma$ (in \kms) is the physical (Doppler) line-width. With
this parameterization, the total flux of each line (in \flux) is
\begin{equation}
  \Phi = \sqrt{2\pi}\,A\,\lambda^*\,\sigma_\ell
       = \sqrt{2\pi}\,A\,\lambda^*\,\frac{\sigma}{c},
  \label{eq:flux_phys}
\end{equation}
which is independent of the DESI wavelength grid and instrumental
resolution.

Comparing this profile to a DESI spectrum requires convolving it with
the resolution matrix $\mathbf{R}$ (\S\ref{ssec:spectra}). Unlike the
stellar continuum, however, which varies smoothly enough across a
pixel that sampling the continuum model at pixel centers before
applying $\mathbf{R}$ preserves its flux (\S\ref{sssec:fitting}), it
is possible for an emission line to be narrower than a single
0.8~\AA{} DESI pixel.\footnote{DESI pixels range in size between
$\approx66$ and $\approx25$~\kms{} over the observed-frame spectral
range, $3600-9800$~\AA.} Low-mass, low-metallicity star-forming
galaxies, in particular, can exhibit strong nebular line-emission with
intrinsic line-widths of tens of \kms{} or less---narrower than a
pixel over most of the observed wavelength range
(\citealt{scholte26a}, J.~Moustakas et~al., in preparation).

The obvious forward-modeling analog, $\mathbf{m} \approx
\mathbf{R}\,\boldsymbol{\phi}$, therefore fails: point-sampling
$\phi(\lambda)$ at the pixel centers captures only a small,
pixel-phase-dependent fraction of a narrow line's true area, biasing
the recovered flux low by a few percent up to $>$30\%, depending on
where the line center falls within a pixel (see Appendix~D in
\citealt{guy23a} for additional discussion). In detail, this bias
arises because $\mathbf{R}$'s rows are normalized to conserve flux
only for smooth inputs sampled at pixel centers, with no matching
column normalization to protect sub-pixel structure
\citep{guy23a}. Moreover, we emphasize that this bias is not just a
modeling artifact; simulations show that even a simple
(model-independent) integration of the observed pixels (BOXFLUX; see
\S\ref{sssec:lines}) suffers the same loss for lines narrower than
about one pixel (S.~Koposov 2026, private
communication).\footnote{\url{https://github.com/desihub/fastspecfit/issues/183}}

%Gaussian-pre-convolved discretization,

To contend with this bias, we adopt the deconvolution strategy devised
by S.~Koposov for \texttt{RVSpecFit} \citep{koposov19b, koposov26a}.
Rather than sampling $\phi(\lambda)$ at pixel centers, we pre-convolve
the model with a fixed Gaussian kernel of width $\sigma_G =
0.5$\,\AA---narrow compared to the DESI resolution but wide enough to
regularize the sub-pixel structure of the model---and propagate the
result through a correspondingly deconvolved resolution matrix
\begin{equation}
  \mathbf{W} = \mathbf{R}\,\mathbf{M}^{-1},
  \label{eq:wmatrix}
\end{equation}
where $\mathbf{M}$ is the Gaussian design matrix, $\mathbf{M}_{ij} =
G(\lambda_i - \lambda_j;\,0,\,\sigma_G)$. The predicted line model is then $\mathbf{m} \approx
\mathbf{W}\,\mathbf{F}$, where $\mathbf{F}$ is the
$\sigma_G$-pre-convolved profile evaluated at pixel centers. Because
each $F_i$ is a weighted integral of $\phi(\lambda)$ rather than a
point sample, it captures the correct line area regardless of the
line's width or its phase (position) relative to the pixel grid.

Formally, $F_i = (\phi * G_{\sigma_G})(\lambda_i)$ is the convolution
of the intrinsic profile given by equation~(\ref{eq:intrinsic}) with
the unit-area Gaussian kernel, evaluated at pixel center
$\lambda_i$. Because both $\phi$ and the kernel are exact Gaussians in
log-$\lambda$ space, their convolution is itself Gaussian, with
variances adding in quadrature to give an effective width
\begin{equation}
  \sigma_\mathrm{eff} = \sqrt{\sigma_\ell^2 + \sigma_{G,\ell}^2},
  \label{eq:sigmaeff}
\end{equation}
\noindent where $\sigma_{G,\ell} \equiv \sigma_G/\lambda^*$. In
addition, because the kernel has unit area, the convolution also
conserves the total line flux, given by
equation~(\ref{eq:flux_phys}). However, since the profile has
broadened from $\sigma_\ell$ to $\sigma_\mathrm{eff}$, its peak
amplitude must shrink by the same factor to keep that flux fixed.
The pre-convolved profile of a single line at pixel center $\lambda_i$
therefore becomes
\begin{equation}
  F_i = A\,\frac{\sigma_\ell}{\sigma_\mathrm{eff}}
    \exp\!\left(-\frac{(\ln\lambda_i-\ln\lambda^*)^2}{2\sigma_\mathrm{eff}^2}\right),
  \label{eq:model}
\end{equation}
directly analogous to equation~(\ref{eq:intrinsic}).  The amplitude
factor $\sigma_\ell/\sigma_\mathrm{eff} \leq 1$ accounts for the
broadening introduced by the pre-convolution kernel. And because
$\sigma_\mathrm{eff} \geq \sigma_{G,\ell}$, even a formally unresolved
line ($\sigma_\ell \to 0$) is represented by a profile spanning
several pixels, so its flux is correctly captured by the pixel grid
rather than lost between samples. Collecting the
pre-convolved profiles of all $n_\mathrm{line}$ lines into
$\mathbf{F}$, the pixelized counterpart of
equation~(\ref{eq:emline_sum}) finally becomes
\begin{equation}
  \mathbf{E} = \mathbf{W}\,\mathbf{F}.
  \label{eq:emline_pixel}
\end{equation}

With this pixelized model in hand, we determine the parameters
$\theta_E$ of every line by minimizing
\begin{equation}
  \chi^2(\theta_E) = \sum_{i=1}^{n_\mathrm{pix}} I_i\bigl(r_i - E_i\bigr)^2,
  \label{eq:chi2line}
\end{equation}
using the trust-region reflective algorithm
\citep[\texttt{scipy.optimize.least\_squares};][]{branch99a}, where
$r_i$ is the residual (continuum-subtracted) spectrum, $I_i$ is the
corresponding inverse variance spectrum, $E_i$ is the emission-line
model given by equation~(\ref{eq:emline_pixel}), and $n_\mathrm{pix}$
is the number of pixels in the spectrum.

Because \fastspecfit{} has to fit tens of blended lines
simultaneously\footnote{Fun fact: at $z \approx 0.076$–$0.083$, as
many as 34 of the 46 lines in the nominal \fastspecfit{} line-list
fall simultaneously within the nominal DESI wavelength range
(3600–9824~\AA).}, this minimization involves many free parameters,
and finite-difference gradients would require $\mathcal{O}(3\times
n_\mathrm{line})$ additional forward-model evaluations per
iteration. We instead speed up the minimization by supplying
\texttt{least\_squares} with an analytic Jacobian: closed-form first
derivatives of equation~(\ref{eq:model}) with respect to each free
parameter, reducing this cost to a single forward-model call per
iteration.

First, because $F_i$ is linear in $A$, the amplitude gradient is
simply
\begin{equation}
  \frac{\partial F_i}{\partial A} = \frac{F_i}{A}.
  \label{eq:dF_dA}
\end{equation}
Writing $x_i \equiv \ln\lambda_i - \ln\lambda^*$ for brevity, the
velocity-shift gradient follows from the dependence of
$\ln\lambda^* = \ln[\lambda_\mathrm{rest}(1 + z + \Delta v/c)]$ on
$\Delta v$:
\begin{equation}
  \frac{\partial F_i}{\partial \Delta v}
  =
  \frac{F_i\,x_i}{\sigma_\mathrm{eff}^2\,c\,(1+z+\Delta v/c)}.
  \label{eq:dF_dv}
\end{equation}
And finally, the line-width gradient, written in a form that remains
finite as $\sigma_\ell\to0$, accounts for the dependence of both the
amplitude factor $\sigma_\ell/\sigma_\mathrm{eff}$ and the exponent on
$\sigma_\ell = \sigma/c$:
\begin{equation}
  \frac{\partial F_i}{\partial \sigma} =
  \frac{A}{\sigma_\mathrm{eff}\,c}\,
  \exp\!\left(-\frac{x_i^2}{2\sigma_\mathrm{eff}^2} \right) \, \Bigl[1
    - \frac{\sigma_\ell^2}{\sigma_\mathrm{eff}^2} \Bigl(1 -
    \frac{x_i^2}{\sigma_\mathrm{eff}^2} \Bigr)\Bigr].
  \label{eq:dF_ds}
\end{equation}

The parameters $\theta_E$ recovered from this minimization are the
amplitude, velocity shift, and width of every individual line. In the
next section (\S\ref{sssec:kinematic}) we describe how we share
$\Delta v$ and $\sigma$ across physically motivated kinematic groups
to reduce this parameter count.

\subsubsection{Kinematic Groups and Broad-Line Decomposition}
\label{sssec:kinematic}

\fastspecfit{} reduces the dimensionality of its emission-line fitting
(\S\ref{sssec:emline_model}) by grouping lines into kinematic groups,
within which all members share a common velocity shift $\Delta v$ and
intrinsic velocity width $\sigma$. All kinematic groups, bounds, and
amplitude constraints are loaded at startup from the bundled
configuration file \texttt{data/emline-constraints.yaml}, as
summarized in Table~\ref{table:kinematic}. This parameter file also
defines two important kinematic groups: an emission line model
consisting entirely of narrow lines (\emph{narrow-only}) and a model
with both narrow and broad lines (\emph{narrow+broad}).

\setlength{\tabcolsep}{1cm}
\begin{deluxetable*}{cc}
\tablecaption{Emission-Line Free Parameters and Bounds\label{table:priors}}
\tablewidth{0pt}
\tablehead{
  \colhead{Parameter} &
  \colhead{Bounds}
}
\startdata
\multicolumn{2}{c}{\textbf{Narrow-line model}} \\
$\sigma_{\mathrm{narrow}}$     & 1--750~\kms \\
$(\Delta v)_{\mathrm{narrow}}$ & $\pm500$~\kms \\
\multicolumn{2}{c}{\textbf{Broad-line model}} \\
$\sigma_{\mathrm{narrow}}$     & 1--750~\kms \\
$\sigma_{\mathrm{broad}}$      & 1--10,000~\kms \\
$(\Delta v)_{\mathrm{narrow}}$ & $\pm500$~\kms \\
$(\Delta v)_{\mathrm{broad}}$  & $\pm2500$~\kms \\
\enddata
\tablecomments{Optimizer bounds for the free kinematic parameters in
  the narrow-only and narrow+broad emission-line models
  (\S\ref{sssec:kinematic}). Broad-line UV and AGN permitted lines use
  the broad-line bounds.}
\end{deluxetable*}
\setlength{\tabcolsep}{6pt}  % restore default


%We first define these groups and the two competing kinematic models,
%then describe how we select between them, and finally an optional
%relaxation pass that refines the winning model's kinematics.

The narrow-only model ties all forbidden lines to one group anchored
to \oiiilam{} and a second group consisting of all \ion{H}{1} and
\ion{He}{1} lines anchored to \ha. ``Anchored'' in this context means
that all the lines share the same $\Delta v$ and $\sigma$.

Meanwhile, the narrow+broad model consists of three groups: one group
consisting of all...forbidden...instead collapses every narrow line
into a single group anchored on narrow \ha, keeps \oiiidoublet{} as
its own group, and adds a third group for broad Balmer emission
anchored on broad \ha.

Narrow-line bounds are $1 < \sigma < 750$~\kms{} and $|\Delta v| <
500$~\kms; broad-line bounds are $1 < \sigma < 10{,}000$~\kms{} and
$|\Delta v| < 2{,}500$~\kms. Within these groups, most doublets---for
example, [\ion{O}{3}]~$\lambda$4959\,/\,$\lambda$5007 and
[\ion{N}{2}]~$\lambda$6548\,/\,$\lambda$6584---carry a fixed amplitude
ratio set by atomic physics, while a few density-sensitive doublets,
such as [\ion{O}{2}]~$\lambda$3726\,/\,$\lambda$3729 and
[\ion{S}{2}]~$\lambda$6716\,/\,$\lambda$6731, instead carry a free
ratio that encodes the electron density. Table~\ref{table:kinematic}
lists every group's full membership, anchor, bounds, and amplitude
constraint for both profile models.

Eight UV and broad-AGN permitted lines carry independently free
kinematics in both fitting models, and two doublet pairs
(\ion{Si}{3}]~$\lambda$1892/\ion{C}{3}]~$\lambda$1908 and
    \mgiidoublet) act as global kinematic anchors, the latter with a
    free amplitude ratio; Table~\ref{table:kinematic} lists these
    lines and bounds in full.

\fastspecfit{} fits every spectrum with both the narrow-only and
narrow+broad models simultaneously, then selects the preferred solution
based on four conditions that must all be satisfied. First, the
reduction in $\chi^2$ over the pixels spanning the broad Balmer lines
must exceed the corresponding reduction in degrees of freedom,
\begin{equation}
  \Delta\chi^2_\mathrm{Balmer} \;>\; \Delta N_\mathrm{dof,\,Balmer},
  \label{eq:broadchi2}
\end{equation}
where $\Delta\chi^2 = \chi^2_\mathrm{narrow} - \chi^2_\mathrm{broad}$
and $\Delta N_\mathrm{dof}$ counts the additional free parameters
introduced by the broad-Balmer group, evaluated only over pixels within
$\pm N_\sigma$ of the non-helium broad Balmer lines. Second, the fitted
velocity dispersion of the broad \ha{] component must exceed
$\sigma_\mathrm{broad}(\ha) > 250$~\kms, excluding spurious broad fits
driven by noise or continuum residuals. Third, the broad component must
be wider than the narrow component,
$\sigma_\mathrm{broad}(\ha) > \sigma_\mathrm{narrow}(\ha)$, ensuring a
physical hierarchy between the two kinematic components. Fourth, at
least one of the two brightest broad Balmer lines in the spectral range
must satisfy $S/N_\mathrm{broad} > 2.5$, where we estimate the noise
robustly from the interquartile range of the fit residuals in the
Balmer-line pixels after subtracting the full broad model. If any
condition fails, we retain the narrow-only result. We record the
pixel-level $\chi^2$ comparison in the output columns
\texttt{DELTA\_LINECHI2} and \texttt{DELTA\_LINENDOF}.

Once a model is selected, \fastspecfit{} runs an optional final
relaxation pass that frees individual velocity shifts for the
narrow-only groups flagged in Table~\ref{table:kinematic}; the
three-group narrow+broad model already has enough kinematic freedom to
absorb residual complexity through its broad Balmer component, and
relaxing it further would risk over-fitting. For the narrow-only
model, we free each narrow line's velocity shift within a
$\pm100$~\kms{] window around its primary best-fit value, while line
widths remain tied within each kinematic group; this accommodates
galaxies with forbidden-line outflows or small wavelength-calibration
residuals, where individual lines can settle at slightly different
velocities than their group's anchor. Doublet-tied pairs (e.g.,
[\ion{N}{2}]~$\lambda\lambda$6548,6584) retain a shared velocity shift
during this relaxation, since releasing transitions from the same
upper level independently would allow unphysical offsets between
them. We adopt the relaxed solution only when $\Delta\chi^2 > \Delta
N_\mathrm{dof}$; otherwise we retain the primary result. We record
this comparison in the output columns \texttt{DELTA\_KINECHI2} and
\texttt{DELTA\_KINENDOF}.

\subsubsection{Summary of Measured Emission-Line Properties}
\label{sssec:lines}

For every line in our list (\S\ref{ssec:emlinelist}), whether formally
detected or not, we tabulate a standard set of derived quantities in
the output catalog; Table~\ref{table:linesummary} summarizes their
symbols, output-column names, and units.

Two amplitude-like quantities are worth distinguishing explicitly. The
model amplitude $A$ is the raw, pre-convolution Gaussian parameter
from equation~(\ref{eq:model}). The observed amplitude $A_\mathrm{obs}$
is instead the peak flux density of the fully pixelized,
resolution-matrix-convolved model $\mathbf{E}=\mathbf{W}\mathbf{F}$
(equation~\ref{eq:emline_pixel}) as it actually appears on the DESI
pixel grid. For lines narrower than a pixel, the pre-convolution
regularization of \S\ref{sssec:emline_model} means $A$ can
substantially exceed $A_\mathrm{obs}$; $A_\mathrm{obs}$, not $A$, is
what one would actually measure by eye in the spectrum.

We likewise report two independent flux estimators. The
Gaussian-integrated flux $\Phi$ is computed analytically from the
fitted $(A,\lambda^*,\sigma)$ via equation~(\ref{eq:flux_phys}),
independent of the pixel grid or resolution matrix. The
boxcar-integrated flux $\Phi_\mathrm{box}$ is instead summed directly
from the continuum-subtracted spectrum,
\begin{equation}
  \Phi_\mathrm{box} = \sum_{i\,\in\,\mathrm{patch}} r_i\,\Delta\lambda_i,
  \label{eq:boxflux}
\end{equation}
over the unmasked pixels within $\pm3\sigma$ of the line center
(requiring at least 7 pixels, and falling back to a default width when
a line was not fit). This is the
same pixel-integrated estimator discussed in
\S\ref{sssec:emline_model}, and it shares the same sub-pixel bias as a
naive point-sampled model for lines narrower than a DESI pixel.

We estimate the local continuum flux density $F_c$ directly from the
data rather than from the smooth continuum model $C(\lambda)$: we take
the sigma-clipped mean of the observed spectrum, with the
best-fitting emission-line model subtracted, over two sideband windows
spanning 3--10$\sigma$ on either side of the line center (again
requiring at least 7 pixels per side). This local, data-driven
estimate is more robust to continuum-model mismatch at the line's
exact wavelength than simply interpolating $C(\lambda)$ would be. The
rest-frame equivalent width then follows as
\begin{equation}
  \mathrm{EW} = \frac{\Phi}{F_c\,(1+z)}.
  \label{eq:ew}
\end{equation}

When a line's amplitude cannot be reliably measured---either because
it was dropped from the fit or its local pixel window is too heavily
masked---we still report a $1\sigma$ upper limit on its flux by
substituting the continuum noise floor for the amplitude in
equation~(\ref{eq:flux_phys}),
\begin{equation}
  \Phi_\mathrm{lim} = \sqrt{2\pi}\,\sigma_\ell\,\lambda^*\,I_{F_c}^{-1/2},
  \label{eq:fluxlimit}
\end{equation}
where $I_{F_c}$ is the inverse variance of $F_c$, and for dropped
lines we substitute a default width for $\sigma$ (75~\kms{} for narrow
lines, 1200~\kms{} for broad Balmer lines) so that a limit can still be
reported. We compute $\Phi_\mathrm{lim}$ for essentially every line
with a valid continuum measurement, independent of whether the line
itself is significantly detected, so that users can define their own
significance threshold---e.g., comparing $\Phi$ to $\Phi_\mathrm{lim}$,
or $A_\mathrm{obs}\sqrt{I_{A_\mathrm{obs}}}$ to a chosen S/N
cut---rather than relying on a single cut we impose. Finally, we
record the pixel-level $\chi^2_\mathrm{line}$ of the fit and the
number of unmasked pixels $n_\mathrm{pix}^\mathrm{line}$ contributing
to it; \S\ref{ssec:uncertainties} describes how the companion
inverse-variance columns flagged in Table~\ref{table:linesummary} are
estimated via our per-object Monte Carlo procedure.

\begin{deluxetable}{lccc}
\tablecaption{Per-Line Measurements\label{table:linesummary}}
\tablewidth{0pt}
\tablehead{
  \colhead{Symbol} &
  \colhead{Column\tablenotemark{a}} &
  \colhead{Units} &
  \colhead{Inverse variance?}
}
\startdata
$A$                       & \texttt{MODELAMP}    & \fluxdensity & \nodata \\
$A_\mathrm{obs}$          & \texttt{AMP}         & \fluxdensity & Y \\
$\Phi$                    & \texttt{FLUX}        & \flux        & Y \\
$\Phi_\mathrm{box}$       & \texttt{BOXFLUX}     & \flux        & Y \\
$\Delta v$                & \texttt{VSHIFT}      & \kms         & Y \\
$\sigma$                  & \texttt{SIGMA}       & \kms         & Y \\
$F_c$                     & \texttt{CONT}        & \fluxdensity & Y \\
EW                        & \texttt{EW}          & \AA          & Y \\
$\Phi_\mathrm{lim}$       & \texttt{FLUX\_LIMIT} & \flux        & \nodata \\
$\chi^2_\mathrm{line}$    & \texttt{CHI2}        & \nodata      & \nodata \\
$n_\mathrm{pix}^\mathrm{line}$ & \texttt{NPIX}   & \nodata      & \nodata \\
\enddata
\tablenotetext{a}{Each column is prefixed by the line name (e.g.,
\texttt{OIII\_5007\_FLUX}); Table~\ref{table:emlines} lists the full
set of modeled lines.}
\tablecomments{Columns marked ``Y'' have a companion \texttt{\_IVAR}
column storing the inverse variance, estimated via the same
per-object Monte Carlo procedure described in \S\ref{ssec:uncertainties}.
$A$, $\Delta v$, and $\sigma$ are defined in \S\ref{sssec:emline_model};
$A_\mathrm{obs}$, $\Phi$, $\Phi_\mathrm{box}$, $F_c$, EW, $\Phi_\mathrm{lim}$,
$\chi^2_\mathrm{line}$, and $n_\mathrm{pix}^\mathrm{line}$ are defined in
\S\ref{sssec:lines}.}
\end{deluxetable}


\subsection{Spectrophotometric and Physical Properties}
\label{ssec:quantities}

Beyond the best-fitting stellar continuum model
$C(\lambda;\theta_C)$ itself (\S\ref{sssec:fitting}), \fastspecfit{}
derives a broad suite of physical and spectrophotometric quantities
and reports them in the output catalog
(\S\ref{sec:appendix_datamodel}). These fall into four groups:
stellar-population parameters---stellar mass, SFR, and age---computed
directly from the SFH template coefficients $\mathbf{c}$; K-corrections
and rest-frame absolute magnitudes, computed by comparing synthesized
model photometry in the observed and rest frames; monochromatic
luminosities and continuum flux densities at fixed rest-frame
wavelengths, measured directly from $C(\lambda)$; and the $D_n(4000)$
spectral break, a purely empirical index that we measure independently
from both the data and the model. We describe each of these in turn.

Recall from \S\ref{sssec:fitting} that each template coefficient $c_n$
carries units of \msun{} of stars formed in age bin $n$
(equation~\ref{eq:fnorm}), rescaled by the mass normalization
$\mathcal{M}_0$ introduced in equation~(\ref{eq:zfactor}). Converting
these coefficients into a present-day stellar mass requires correcting
for the mass each population has since returned to the ISM via stellar
winds and supernovae; we take these surviving-mass fractions
$f_{\star,n}\in(0,1]$ directly from FSPS and compute
\begin{equation}
  \Mstar = \mathcal{M}_0\sum_n c_n\,f_{\star,n}.
  \label{eq:mstar}
\end{equation}
We caution that non-parametric SFH models can yield stellar masses
that are systematically higher---by $\sim$0.2--0.3~dex on
average---than parametric estimates, because a flexible SFH can
accommodate more mass in old, low-luminosity stellar populations
\citep{leja19a,leja20a}. As noted in \S\ref{ssec:templates}, our
current SPS template grid is fixed at solar metallicity, so none of
the quantities derived here---or their reported uncertainties---carry
any information about metallicity.

We estimate the star-formation rate (SFR) by summing the mass formed
within the most recent 100~Myr, weighting each age bin by the fraction
of its width that falls within that window (for our current age
binning, \S\ref{ssec:templates}, this reduces to the two youngest
bins, 0--30~Myr and 30--100~Myr, each fully included), and dividing by
100~Myr,
\begin{equation}
  \mathrm{SFR} = \frac{\mathcal{M}_0}{100\,\mathrm{Myr}}\sum_n c_n\,w_n,
  \label{eq:sfr}
\end{equation}
where $w_n\in[0,1]$ is the fraction of age bin $n$ within 100~Myr of
observation. Although \fastspecfit{} measures both the \ha{}
luminosity and the rest-frame UV continuum---each of which serves as
an independent SFR diagnostic---we do not currently convert either
quantity to a SFR within the code.

Finally, we compute a mean stellar age as the coefficient-weighted
average of the age-bin midpoints $t_n$,
\begin{equation}
  \langle t_\star\rangle = \frac{\sum_n c_n\,t_n}{\sum_n c_n}.
  \label{eq:age}
\end{equation}
Because the coefficients $c_n$ represent formed stellar mass rather
than surviving mass or luminosity, this is neither a strictly
light-weighted nor a strictly (surviving-)mass-weighted age, but falls
between the two.

We compute K-corrections and rest-frame absolute magnitudes by
comparing synthesized photometry from the best-fitting continuum model
$C(\lambda)$ in the observed and rest frames. The output rest-frame
bandpasses and their band-shifts are loaded at startup from the
bundled configuration file \texttt{data/legacysurvey-dr9.yaml}, and
span DECam $grz$ (band-shifted to $z=1$), Johnson/Cousins $UBV$ and
2MASS $J$ (band-shifted to $z=0$), and SDSS $ugriz$ and WISE $W1$
(band-shifted to $z=0.1$, following \citealt{blanton07a}); the
corresponding output columns are documented in full in
Table~\ref{table:datamodel} (Appendix~\ref{sec:appendix_datamodel}),
e.g., the rest-frame SDSS $g$-band magnitude band-shifted to $z=0.1$
is recorded as \texttt{ABSMAG01\_SDSS\_G}. For each band we identify
the observed-frame Legacy Survey or WISE bandpass whose rest-frame
effective wavelength most closely matches the target band, restricted
to bandpasses detected at $\mathrm{S/N}>2$; this minimizes the
(typically small) K-correction required. The absolute magnitude then
follows the standard relation \citep{hogg02a},
\begin{equation}
  M = m_\mathrm{obs} - \mathrm{DM}(z) - K,
  \label{eq:kcorr}
\end{equation}
where $m_\mathrm{obs}$ is the observed-frame photometry,
$\mathrm{DM}(z)$ is the distance modulus, and
$K = 2.5\log_{10}(f_\mathrm{rest}^\mathrm{synth}/f_\mathrm{obs}^\mathrm{synth})$
is the K-correction, with $f_\mathrm{rest}^\mathrm{synth}$ and
$f_\mathrm{obs}^\mathrm{synth}$ the rest- and observed-frame fluxes
synthesized from $C(\lambda)$ in the target and nearest-matching
bandpasses, respectively. When no bandpass satisfies the S/N
threshold, we instead report the absolute magnitude synthesized purely
from the model, and set its inverse variance to zero.

To construct $f_\mathrm{rest}^\mathrm{synth}$ in a band-shifted system
\citep{blanton07a}, we blueshift the filter response itself rather
than the model, replacing the native response $R_Q(\lambda)$ of
bandpass $Q$ with
\begin{equation}
  R_Q^{(q)}(\lambda) \equiv R_Q\bigl((1+q)\,\lambda\bigr),
  \label{eq:bandshift}
\end{equation}
so that convolving $R_Q^{(q)}$ with the rest-frame SED reproduces what
$R_Q$ would measure on the same galaxy at redshift $q$; $q=0$ recovers
the native, unshifted bandpass. Table~\ref{table:absmag-bands} lists
the band-shift $q$ adopted for each of our output bandpasses, together
with their native filter curves and effective wavelengths.

\begin{deluxetable}{llcc}
\tablecaption{Rest-Frame Bandpasses for Absolute Magnitudes and K-Corrections\label{table:absmag-bands}}
\tablewidth{0pt}
\tablehead{
\colhead{Band} &
\colhead{Filter Curve} &
\colhead{$\lambda_\mathrm{eff}$\tablenotemark{a}} &
\colhead{Band-shift $z_0$}
}
\startdata
DECam $g$    & \texttt{decam2014-g} & 4890  & 1.0 \\
DECam $r$    & \texttt{decam2014-r} & 6470  & 1.0 \\
DECam $z$    & \texttt{decam2014-z} & 9196  & 1.0 \\
Johnson $U$  & \texttt{bessell-U}   & 3618  & 0.0 \\
Johnson $B$  & \texttt{bessell-B}   & 4442  & 0.0 \\
Johnson $V$  & \texttt{bessell-V}   & 5536  & 0.0 \\
2MASS $J$    & \texttt{twomass-J}   & 12408 & 0.0 \\
SDSS $u$     & \texttt{sdss2010-u}  & 3614  & 0.1 \\
SDSS $g$     & \texttt{sdss2010-g}  & 4749  & 0.1 \\
SDSS $r$     & \texttt{sdss2010-r}  & 6206  & 0.1 \\
SDSS $i$     & \texttt{sdss2010-i}  & 7526  & 0.1 \\
SDSS $z$     & \texttt{sdss2010-z}  & 8947  & 0.1 \\
WISE $W1$    & \texttt{wise2010-W1} & 34003 & 0.1 \\
\enddata
\tablecomments{Default rest-frame bandpasses, band-shifts $z_0$, and
  \texttt{speclite} filter curves used to infer K-corrections and
  absolute magnitudes (\S\ref{ssec:quantities}), as configured in the
  bundled \texttt{data/legacysurvey-dr9.yaml} parameter file.}
\tablenotetext{a}{Effective filter wavelengths in \AA{} are reported
  for the unshifted filter curve.}
\end{deluxetable}


We additionally measure rest-frame monochromatic luminosities and
observed-frame continuum flux densities directly from $C(\lambda)$,
using a robust (sigma-clipped) linear fit to the continuum within
$\pm100$~\AA{} of each reference wavelength---in some cases
(Ly$\alpha$, \hb, \ha) excluding the inner $\pm50$~\AA{} to avoid bias
from strong, broad emission-line wings. We report luminosities at
1450, 1700, 3000, and 5100~\AA{} in solar units ($\lambda L_\lambda$),
and at 1500 and 2800~\AA{} in $\mathrm{erg\,s^{-1}\,Hz^{-1}}$
($L_\nu$), the latter chosen to facilitate UV-continuum SFR estimates
(e.g., \citealt{kennicutt98a}). We similarly report the observed-frame
continuum flux density at the rest-frame centers of Ly$\alpha$, \oii,
\hb, \oiiilam, and \ha, which provide simple, spectral-model-independent
continuum-level diagnostics at the locations of these key nebular
emission lines.

Finally, we measure the narrow definition of the 4000~\AA{} break,
$D_n(4000)$, the ratio of the mean flux density in the
4000--4100~\AA{} and 3850--3950~\AA{} rest-frame windows
\citep{balogh99a}, a robust, model-independent tracer of
light-weighted stellar age that is largely insensitive to reddening.
We report three flavors of this index: \texttt{DN4000\_OBS}, measured
directly from the observed spectrum, without any correction for
emission-line contamination; \texttt{DN4000\_MODEL}, measured from the
best-fitting stellar continuum model alone (line- and noise-free, but
dependent on the SPS templates); and our fiducial \texttt{DN4000},
measured from the observed spectrum after subtracting the
best-fitting emission-line model, combining the empirical fidelity of
the data with a correction for line contamination within the break
windows.

\subsection{Parameter Uncertainties}
\label{ssec:uncertainties}

Because \fastspecfit{}'s forward model chains together several
nonlinear stages---an initial grid search for \vdisp{}
(\S\ref{sssec:vdisp}), the VARPRO solve for \tauv{} and $\mathbf{c}$
(\S\ref{sssec:fitting}), and the nonlinear emission-line kinematics
(\S\ref{sssec:kinematic})---propagating uncertainties analytically
through every downstream quantity (stellar mass, SFR, K-corrections,
equivalent widths, etc.) is impractical. We instead adopt a Monte
Carlo approach: by perturbing the input data and refitting, each
nonlinear stage of the pipeline propagates its own uncertainty
naturally, without our having to derive a bespoke error-propagation
formula for every derived quantity.

For each object we generate $N_\mathrm{monte}=50$ independent
realizations of the observed spectrum and photometry by adding
Gaussian noise consistent with the pipeline's inverse-variance
estimates,
\begin{equation}
  d_i^{(m)} \sim \mathcal{N}\bigl(d_i,\,I_i^{-1}\bigr),
  \label{eq:mc_draw}
\end{equation}
for realizations $m=1,\ldots,N_\mathrm{monte}$ (and similarly for the
photometry). These draws are generated once, from a reproducible
per-object seed recorded in the \texttt{SEED} output column---itself
derived deterministically from a single global run seed---and are
reused, unchanged, across every subsequent fitting stage for that
object: the same $N_\mathrm{monte}$ realizations flow through the
velocity-dispersion, continuum, and emission-line fits described
below. The inverse variance of any fitted or derived quantity $\theta$
then follows from the scatter across realizations,
\begin{equation}
  \sigma_\theta^{-2}
  = \bigl[\mathrm{Var}(\theta^{(1)},\ldots,\theta^{(N_\mathrm{monte})})\bigr]^{-1}.
  \label{eq:mc_unc}
\end{equation}

On the continuum side, the Monte Carlo realizations enter in two
places. First, while \vdisp{} is being determined (\S\ref{sssec:vdisp}),
we refit \tauv, \vdisp, and $\mathbf{c}$ jointly on each realization to
obtain \texttt{VDISP\_IVAR}. Second, with \vdisp{} fixed at its adopted
value, we repeat the full VARPRO solve---the outer \tauv{} minimization
and inner NNLS solve for $\mathbf{c}$---on each realization, yielding
$N_\mathrm{monte}$ independent estimates $\tauv^{(m)}$ and
$\mathbf{c}^{(m)}$ and directly giving \texttt{TAUV\_IVAR}. Because
every stellar-population and photometric quantity in
\S\ref{ssec:quantities} is a fixed function of $\mathbf{c}$
(equations~\ref{eq:mstar}--\ref{eq:bandshift}), we simply re-evaluate
those relations on each $\mathbf{c}^{(m)}$ and propagate the resulting
scatter to \texttt{LOGMSTAR\_IVAR}, \texttt{SFR\_IVAR},
\texttt{AGE\_IVAR}, \texttt{DN4000\_MODEL\_IVAR}, the
\texttt{ABSMAG*\_SYNTH\_IVAR} and rest-frame luminosity/continuum-flux
inverse variances---no additional nonlinear refitting is required.

The emission-line Monte Carlo is deliberately cheaper. Rather than
repeating the full kinematic model-selection procedure
(\S\ref{sssec:kinematic})---which chooses between narrow-only and
narrow+broad models and optionally relaxes velocity-width
bounds---for every realization, we fix the kinematic model and bounds
at their fiducial, data-preferred values and refit only the line
parameters, warm-started at the fiducial best-fit solution with
loosened convergence tolerances. This is a good approximation because
each realization is a small perturbation of the data rather than an
independent object, so the preferred kinematic model rarely changes
between realizations. We build each realization's continuum-subtracted
spectrum from the corresponding continuum-fitting realization, $d^{(m)}
- C^{(m)}(\lambda) - S(\lambda)$, holding the smooth continuum
correction $S(\lambda)$ (\S\ref{sssec:smooth}) fixed at its fiducial
value across all realizations rather than recomputing it per draw.
This procedure yields Monte Carlo uncertainties for every fitted line
parameter---amplitude, flux, equivalent width, velocity shift and
width, non-parametric moments---and for the doublet ratios.

\subsection{Additional Fitting Modes}
\label{sec:modes}

\fastspecfit{} operates in two modes. \fastphot{} fits broadband SEDs
to Legacy Survey photometry with no spectroscopic data. \fastspec{}
jointly fits spectra from all three DESI cameras (B, R, Z) and
broadband photometry simultaneously. Both modes share the same
template library, dust model, and physical output quantities; they
differ only in the data that enter the fit and in whether the stellar
velocity dispersion is a free parameter.

In \fastphot{} mode we do not fit a velocity dispersion: we
convolve all templates with ${\vdisp}^\mathrm{kern}_\mathrm{nom}$ and
report $\vdisp = 150$~\kms{} with \texttt{VDISP\_IVAR}~$= 0$.

A handful of options adapt \fastspecfit{} to non-default inputs. The
\texttt{--targetids} and \texttt{--input-redshifts} options fit a
specified subset of targets at externally supplied redshifts (e.g.,
from visual inspection) instead of the pipeline's Redrock or
QuasarNet value. \texttt{--ignore-photometry} drops the broadband
photometry from the fit; because the aperture correction
(\S\ref{sssec:apercorr}) requires comparing the model to the observed
photometry, this also disables it, leaving the spectrum on its
native, uncalibrated flux scale---we do not use this option in the
DR2 production runs. And \texttt{--fphotofile} substitutes an
alternate photometric-catalog configuration file for the bundled
\texttt{data/legacysurvey-dr9.yaml} (\S\ref{ssec:quantities}), letting
\fastspecfit{} fit photometry from surveys other than the Legacy
Surveys. The DESI coaddition type---per-exposure, per-night,
cumulative (multi-night), or per-HEALPix---is detected automatically
from the input Redrock file and requires no separate configuration.

Finally, a third entry point, \texttt{stackfit}, fits arbitrary
user-supplied stacked spectra (e.g., binned by stellar mass and
redshift) rather than individual DESI targets; since such stacks lack
associated broadband photometry, \texttt{stackfit} runs with
\texttt{--ignore-photometry} enabled by default and fits the spectrum
on its own internal flux scale.

\todo{Note that the user can easily define their own custom
  line-list.}

%% ============================================================
\section{Illustrative Examples}
\label{sec:casestudy}
%% ============================================================

We illustrate the behavior of \fastspecfit{} across the full range of DESI
extragalactic targets with three representative fits spanning nearly two decades
in redshift: a low-redshift BGS galaxy with prominent AGN broad-line emission
(\S\ref{ssec:example_bgs}); an intermediate-redshift ELG at $z \approx 0.8$,
where several emission lines besides the \oii{} doublet fall within the DESI
spectral window (\S\ref{ssec:example_elg}); and a high-redshift QSO
(\S\ref{ssec:example_qso}). Together these examples highlight both the dynamic
range of the code and its known limitations---most notably that the SPS
templates are entirely stellar, so users should interpret continuum-derived
quantities for QSOs and broad-line AGN with appropriate caution.

\subsection{Low-Redshift BGS Galaxy}
\label{ssec:example_bgs}

\begin{figure*}
\includegraphics[width=6.5in]{figures/example-bgs.pdf}
\caption{\fastspecfit{} output for a low-redshift BGS galaxy at
  $z=0.1201$ (TARGETID~39627813136892057; main survey, bright program,
  healpix~31152).  Top: broadband SED, with the best-fit stellar
  continuum model (gray curve), synthesized photometry (open
  diamonds), and observed Legacy Surveys $grz$ and WISE $W1$--$W4$
  photometry with $1\sigma$ error bars (filled orange circles); the
  top axis gives the rest-frame wavelength, and the gray bar below the
  photometry marks the observed DESI wavelength range and the
  $2.43\times$ aperture-correction factor. The inset in the lower
  right shows the Legacy Survey $grz$ image with the
  $1\farcs5$-diameter DESI fiber (solid) and a $10\arcsec$ reference
  aperture (dashed) overlaid. Middle: full observed-frame DESI
  spectrum (light colors, one per camera) with the best-fit stellar
  continuum (black), continuum-plus-smooth-continuum correction
  (dashed gray), and total spectrophotometric model including emission
  lines (bold, per-camera colors); light gray labels mark every
  in-range emission line included in the fit.  Bottom: zooms on four
  representative line complexes---\oiidoublet{}, \hb+\oiii{},
  \ha+\nii{}, and \sii{}~$\lambda\lambda6716,31$---showing the
  continuum- and smooth-continuum-subtracted data, individual best-fit
  Gaussian components (thin), and the total emission-line model
  (bold).\label{fig:example_bgs}}
\end{figure*}

Figure~\ref{fig:example_bgs} shows the \fastspecfit{} output for a
low-redshift BGS galaxy at $z=0.1201$ observed with the DESI fiber
centered on the nucleus. The top panel shows the broadband SED, with
the Legacy Survey $grz$ image inset in the lower right showing the
$1\farcs5$-diameter fiber footprint; the middle panel displays the
full observed DESI spectrum across all three cameras together with the
best-fit spectrophotometric model; and the bottom row zooms in on four
representative emission-line complexes. This object provides a clean
end-to-end illustration of the complete fitting sequence described in
\S\ref{sec:modeling}: the high signal-to-noise stellar continuum
($S/N_b=6.4$, $S/N_r=13.8$, $S/N_z=19.3$) reveals prominent absorption
features, the emission-line spectrum shows both narrow forbidden lines
and broad Balmer emission indicative of a type~1 AGN, and strong WISE
detections anchor the infrared SED.

We walk through each stage of the fitting sequence in turn. The line masker
(\S\ref{ssec:patches}) identifies and excludes the emission-line-affected
pixels before the stellar continuum is fit; the resulting velocity
dispersion, $\vdisp=105$~\kms{} (\S\ref{sssec:vdisp}), and the joint
spectrophotometric continuum fit (\S\ref{sssec:fitting}) together return
$\logmstar=10.19\pm0.02$, $\tauv=1.60\pm0.04$, a light-weighted age of
$4.6\pm0.3$~Gyr, and $\mathrm{SFR}=39\pm2$~\msun{}~yr$^{-1}$. The median
aperture correction, $2.43\times$ (the gray bar beneath the SED in
Figure~\ref{fig:example_bgs}), places the DESI fiber spectrum on the
total-flux scale of the broadband photometry (\S\ref{sssec:apercorr}), and
the smooth continuum correction (\S\ref{sssec:smooth}) then absorbs the
small residual flux-calibration offsets we measure across the three cameras.

The emission-line fit (\S\ref{ssec:emlines}) reveals this object's most
distinctive feature: a broad Balmer component with
$\sigma_\mathrm{broad}(\ha)=880$~\kms{}, well above both the $250$~\kms{}
minimum-width threshold and the narrow component's
$\sigma_\mathrm{narrow}(\ha)=87\pm71$~\kms{}. Detected at $S/N=6.6$ in broad
\ha{} (and only $S/N=1.1$ in broad \hb, still sufficient since only one of
the two brightest broad Balmer lines needs to satisfy the $S/N>2.5$
threshold), the broad component drives $\Delta\chi^2_\mathrm{Balmer}=2751$
against $\Delta N_\mathrm{dof,\,Balmer}=5$ (Equation~\ref{eq:broadchi2}), and
\fastspecfit{} adopts the narrow+broad kinematic model
(\S\ref{sssec:kinematic}). The bottom row of Figure~\ref{fig:example_bgs}
shows the resulting broad wings underlying the \ha+\nii{} and \hb+\oiii{}
complexes alongside the resolved \oiidoublet{} and \sii{}~$\lambda\lambda6716,31$
doublets. The correspondingly high $D_n(4000)=1.27$ (\S\ref{ssec:quantities})
is consistent with the moderately dust-attenuated ($\tauv=1.60$),
intermediate-age stellar population recovered by the continuum fit.

\subsection{Intermediate-Redshift Emission-Line Galaxy}
\label{ssec:example_elg}

\begin{figure*}
\includegraphics[width=6.5in]{figures/example-elg.pdf}
\caption{\fastspecfit{} output for an intermediate-redshift ELG at
  $z=0.9449$ (TARGETID~39627663165362386; main survey, dark program,
  healpix~29973). Top: broadband SED, with the best-fit stellar
  continuum model (gray curve), synthesized photometry (open
  diamonds), and observed Legacy Surveys $grz$ and WISE $W1$--$W4$
  photometry shown either as detections with $1\sigma$ error bars
  (filled orange circles) or as upper limits (open circles with
  downward carets); the top axis gives the rest-frame wavelength, and
  the gray bar below the photometry marks the observed DESI wavelength
  range and the $1.03\times$ aperture-correction factor. The inset in
  the lower right shows the Legacy Survey $grz$ image with the
  $1\farcs5$-diameter DESI fiber (solid) and a $10\arcsec$ reference
  aperture (dashed) overlaid. Middle: full observed-frame DESI
  spectrum (light colors, one per camera) with the best-fit stellar
  continuum (black), continuum-plus-smooth-continuum correction
  (dashed gray), and total spectrophotometric model including emission
  lines (bold, per-camera colors); light gray labels mark every
  in-range emission line included in the fit. Bottom: zooms on four
  representative lines---\oiidoublet{}, H$\gamma$, \hb{}, and
  \oiiidoublet{}---showing the continuum- and
  smooth-continuum-subtracted data, individual best-fit Gaussian
  components (thin), and the total emission-line model
  (bold).\label{fig:example_elg}}
\end{figure*}

Figure~\ref{fig:example_elg} shows an intermediate-redshift ELG at
$z=0.9449$, where the DESI spectral window covers rest-frame
$\sim$1900--5000~\AA{} and brings into view several emission lines beyond
the \oiidoublet{} doublet, including H$\gamma$, \hb{}, and \oiiidoublet{}.
The faint, blue star-forming continuum---median $S/N_b=2.5$, $S/N_r=3.3$,
and $S/N_z=3.7$, roughly a factor of five fainter than the BGS example in
\S\ref{ssec:example_bgs}---tests \fastspecfit{}'s ability to recover
accurate line fluxes and a physically meaningful stellar mass in the
low-continuum regime typical of ELGs.

Even at this low continuum signal-to-noise, the joint spectrophotometric
continuum fit (\S\ref{sssec:fitting}) converges to a young, low-mass, dusty
stellar population: $\logmstar=9.51\pm0.19$, $\tauv=0.37\pm0.02$, a
light-weighted age of only $0.06$~Gyr (poorly constrained, consistent with
zero within its $3.9$~Gyr uncertainty), and
$\mathrm{SFR}=39.2\pm0.9$~\msun{}~yr$^{-1}$, giving the correspondingly low
$D_n(4000)=1.08$
(\S\ref{ssec:quantities}) visible in the middle panel. The fitted stellar
velocity dispersion, $\vdisp=85$~\kms{} (\S\ref{sssec:vdisp}), is
consistent with the fitted narrow emission-line width,
$\sigma_\mathrm{narrow}=101\pm5$~\kms{}. The median aperture correction is
only $1.03\times$ (the gray bar beneath the SED), reflecting how little
flux this compact, low-surface-brightness target loses outside the DESI
fiber, and the smooth continuum correction (\S\ref{sssec:smooth}) absorbs
percent-level residual flux-calibration offsets across all three cameras
($-4.7\%$, $+7.4\%$, and $-4.4\%$ in $b$, $r$, and $z$, respectively).

With \ha{} redshifted beyond the DESI window at this redshift,
\fastspecfit{} cannot evaluate the broad-line selection criteria that are
anchored on \ha{} (\S\ref{sssec:kinematic}) and correctly retains the
narrow-only kinematic model. The bottom row of Figure~\ref{fig:example_elg}
shows the resulting narrow-line fits, including the resolved \oiidoublet{}
doublet, whose free amplitude ratio encodes the electron density of the
star-forming gas (\S\ref{sssec:kinematic}), and the resolved \oiiidoublet{}
doublet with its fixed $\sim1{:}3$ amplitude ratio set by atomic physics.
The recovered stellar mass, star-formation rate, and $D_n(4000)$ are all
consistent with a young, low-mass, actively star-forming galaxy.

\subsection{High-Redshift Quasar}
\label{ssec:example_qso}

\begin{figure*}
\includegraphics[width=6.5in]{figures/example-qso.pdf}
\caption{\fastspecfit{} output for a high-redshift QSO at $z=2.0731$
  (TARGETID~39627657142340435; main survey, dark program,
  healpix~29972).  Top: broadband SED, with the best-fit stellar
  continuum model (gray curve), synthesized photometry (open
  diamonds), and observed Legacy Surveys $grz$ and WISE $W1$--$W4$
  photometry shown either as detections with $1\sigma$ error bars
  (filled orange circles) or as upper limits (open circles with
  downward carets); the top axis gives the rest-frame wavelength, and
  the gray bar below the photometry marks the observed DESI wavelength
  range and the $1.14\times$ aperture-correction factor. The inset in
  the lower right shows the Legacy Survey $grz$ image with the
  $1\farcs5$-diameter DESI fiber (solid) and a $10\arcsec$ reference
  aperture (dashed) overlaid. Middle: full observed-frame DESI
  spectrum (light colors, one per camera) with the best-fit stellar
  continuum (black), continuum-plus-smooth-continuum correction
  (dashed gray), and total spectrophotometric model including emission
  lines (bold, per-camera colors); light gray labels mark every
  in-range emission line included in the fit. Bottom: zooms on four
  broad UV line complexes---Ly$\alpha$+\ion{N}{5}, \ion{C}{4},
  \ion{C}{3}]+\ion{Si}{3}]+ \ion{Al}{3}, and \mgii---showing the data,
  the individual best-fit Gaussian components (thin), and the total
  emission-line model (bold).\label{fig:example_qso}}
\end{figure*}

Figure~\ref{fig:example_qso} shows a QSO at $z=2.0731$, where four
broad UV emission-line
complexes---Ly$\alpha$+\ion{N}{5}~$\lambda\lambda1216,1240$,
\ion{C}{4}~$\lambda1549$, \ion{C}{3}]~$\lambda1908$ blended with
  \ion{Si}{3}]~$\lambda1892$ and \ion{Al}{3}~$\lambda1857$, and
    \mgiidoublet---fall within the DESI spectral window. The spectrum
    is bright (median $S/N_b=16.9$, $S/N_r=20.7$, and $S/N_z=21.4$)
    and the top panel shows the smooth, blue power-law continuum
    characteristic of an unobscured, accretion-disk-dominated quasar,
    rising steadily from the optical through the mid-infrared.

This example illustrates both the flexibility and a fundamental limitation
of \fastspecfit{}'s continuum model. The fitting machinery itself makes no
assumption about the target class: the same non-negative linear
combination of stellar population synthesis templates
(\S\ref{sssec:fitting}) that works for the star-forming and quiescent
galaxies in \S\ref{ssec:example_bgs} and \S\ref{ssec:example_elg} runs here
without any special-casing and formally converges to a well-defined
solution. For a QSO, however, the true continuum is not a non-negative sum
of stellar templates: the accretion disk contributes a significant
non-thermal power-law component that our purely stellar template set
cannot represent, so \fastspecfit{} instead absorbs this excess blue flux
into young, dusty, high-mass stellar-population coefficients. The
resulting continuum solution is, in detail, wrong, even though it is
formally well constrained. The fitted parameters for this object make the
point concretely: $\logmstar=11.13\pm0.00$, $\tauv=0.89\pm0.00$, a
light-weighted age pinned at the youngest template age
($0.02\pm0.00$~Gyr), and an implied
$\mathrm{SFR}=1503\pm9$~\msun{}~yr$^{-1}$---a star-formation rate far in
excess of any known star-forming galaxy, with formal uncertainties that
are vanishingly small precisely
because the fit has enough freedom (many templates, a wide dust range) to
lock onto a single, reproducible, but unphysical solution. Continuum-derived
quantities are therefore not reliable for luminous, unobscured QSOs and
broad-line AGN; addressing this properly requires a non-thermal continuum
component and is planned for a future, dedicated \texttt{fastqso} fitting
mode (\S\ref{ssec:improvements}).

The broad emission-line measurements are more directly useful, though
still approximate. With no Balmer lines falling within the DESI window at
this redshift, \fastspecfit{} skips the narrow-versus-broad Balmer
selection (\S\ref{sssec:kinematic}) entirely and instead fits the four UV
line complexes with their own always-broad kinematics, returning line
widths of order $\sigma_\mathrm{UV}\approx2700$~\kms{} and a velocity
offset of order $\Delta v_\mathrm{UV}\approx-640$~\kms{}, each with a
large formal uncertainty ($1500$--$1300$~\kms{}). These large uncertainties
reflect a real limitation of the model: our single-Gaussian line profile is
a poor approximation to the strong, often asymmetric UV/QSO broad lines
visible in the bottom row of Figure~\ref{fig:example_qso}, and especially
to Ly$\alpha$, a resonance line whose radiative transfer through the
intergalactic and circumgalactic medium routinely produces complex,
asymmetric profiles that no single Gaussian can capture. Users requiring
precise UV broad-line kinematics should treat \fastspecfit{}'s
single-component measurements as an effective width and centroid rather
than a physically complete line-profile decomposition.


%% ============================================================
\section{DR2 Value-Added Catalog}
\label{sec:vac}
%% ============================================================

In this section we present the \fastspecfit{} VAC for the 

DESI spectra are processed in spectroscopic productions, or
specprods: fixed sets of input spectra and calibration files reduced
with a tagged version of the DESI pipeline, named alphabetically after
mountains or mountain ranges. A single data release can include more
than one specprod, and a given specprod can in principle be processed
into more than one \fastspecfit{} VAC. The primary specprod for DR2 is
Loa; the catalog we describe in this section is Loa~v2.0, and we refer
to it throughout as the DR2 VAC. This section describes how we produce
that catalog by running \fastspecfit{} on every DESI extragalactic
target: the sample-selection criteria, our correction of a subset of
\texttt{Redrock} QSO redshifts, and the catalog's production and file
organization (\S\ref{ssec:catalog}), recommended quality cuts
(\S\ref{ssec:cuts}), the code's computational performance at DR2 scale
(\S\ref{ssec:performance}), and our validation of \fastspecfit{}
measurements against independent codes and catalogs
(\S\ref{ssec:validation}).

\subsection{Sample Selection and Catalog Production}
\label{ssec:catalog}

We build the DR2 VAC by applying a uniform sample-selection criterion
to every DESI extragalactic target, correcting a subset of
\texttt{Redrock} QSO redshifts, and fitting the result with
\fastspecfit{}.

We model all extragalactic DESI targets that satisfy three selection criteria: a
spectroscopic redshift $z > 10^{-3}$, a non-sky target class
(\texttt{OBJTYPE}~=~\texttt{TGT}), and a valid spectrum
(\texttt{ZWARN}~\&~\texttt{NODATA}~=~0, where the \texttt{NODATA} bit flags
spectra in which all extracted inverse-variance values are zero). No cuts are
imposed on targeting bits or other redshift-warning bits, ensuring that
\fastspecfit{} results are available for the widest possible range of objects
across every DESI target class.

For a small but astrophysically important fraction of quasar targets,
the \texttt{Redrock} spectral classification pipeline assigns an
incorrect redshift.  We correct these using the \texttt{QuasarNet} and
\mgii afterburner catalogs distributed alongside each \texttt{Redrock}
output file. For primary QSO targets, we replace the \texttt{Redrock}
redshift when two conditions are met: \texttt{QuasarNet} independently
classifies the object as a QSO, and the maximum line confidence across
six diagnostic transitions (Ly$\alpha$, \ion{C}{4}, \ion{C}{3}, \mgii,
\hb, and \ha) exceeds 0.99 for DR2 (0.95 for DR1 and the EDR). A
broader criterion is applied to variable WISE QSO secondary targets;
full details are documented in the online VAC documentation. The
original \texttt{Redrock} redshift and \texttt{ZWARN} bitmask are
preserved in the output columns \texttt{Z\_RR} and
\texttt{ZWARN\_RR}. In total, corrected redshifts are applied to
approximately 9,100 targets in the EDR, 74,000 in DR1, and 165,000 in
DR2.

We produce the DR2 VAC (Loa~v2.0) using \fastspecfit{}
version~3.6.1 with the \texttt{ftemplates-chabrier-2.2.0.fits} SPS template
library (\S\ref{sec:appendix_sps}) and \texttt{desitarget}~v2.8.0 for target
selection and masking, running on the NERSC Perlmutter supercomputer. The
catalog encompasses the complete five-year main survey and totals
approximately 38~million targets, of which the main survey---split between
roughly 13~million BGS targets observed under bright conditions and 23~million
LRG, ELG, and QSO targets observed under dark conditions---contributes
36.5~million. Table~\ref{table:sample} summarizes the complete breakdown by
survey and observing program; the corresponding EDR and DR1 breakdowns appear
in Table~\ref{table:sample-edr-dr1} (Appendix~\ref{sec:appendix_vacs}), where
we also describe key differences in template versions and algorithmic
choices relative to DR2.

\begin{deluxetable*}{lrr}
\tablecaption{Summary of the DR2/Loa \fastspecfit{} Value-Added Catalog\label{table:sample}}
\tablewidth{0pt}
\tablehead{
  \colhead{Program} &
  \colhead{\shortstack{Number of\\Objects}} &
  \colhead{\shortstack{Number of Corrected\\QSO Redshifts\tablenotemark{a}}}
}
\startdata
\cutinhead{Main Survey}
backup &     60,372 &        0 \\
bright & 13,185,499 &    1,161 \\
dark   & 23,273,230 &  157,118 \\
total  & 36,519,101 &  158,279 \\
\cutinhead{Special Survey}
backup &        512 &        0 \\
bright &    105,172 &        1 \\
dark   &    261,706 &      232 \\
total  &    367,390 &      233 \\
\cutinhead{Survey Validation\tablenotemark{b}}
backup &      2,825 &       17 \\
bright &    438,456 &       75 \\
dark   &    880,642 &    6,009 \\
other  &     36,769 &      106 \\
total  &  1,358,692 &    6,207 \\
\tableline
Total & 38,245,183 &  164,719 \\
\enddata
\tablenotetext{a}{Number of corrected QSO redshifts (see \S\ref{ssec:catalog}.}
\tablenotetext{b}{Sum of CMX, SV1, SV2, and SV3 surveys.}
\tablecomments{Target counts in the \fastspecfit{} DR2 VAC, organized
  by survey and observing program.}

\end{deluxetable*}


We run \fastspecfit{} on the HEALPix-coadded DESI spectra and organize the
output catalogs identically to how the underlying spectra, redshift catalogs,
and other data products are organized within DESI data releases. Most users
will find it most convenient to work with the merged catalogs, which combine
all per-HEALPix files for a given \texttt{SURVEY} and \texttt{PROGRAM} into a
single file:

\medskip
\noindent\texttt{VACDIR/catalogs/\{fastspec,fastphot\}-SURVEY-PROGRAM.fits}
\medskip

\noindent where \texttt{VACDIR} is the top-level VAC directory. We
make all catalogs publicly accessible via
\url{https://data.desi.lbl.gov}.
Appendix~\ref{sec:appendix_datamodel} describes the per-HEALPix file
organization, the \texttt{nside}~=~1 subdivision used for the largest
survey--program combinations, and the complete output data model.

\subsection{Recommended Quality Cuts}
\label{ssec:cuts}

We recommend science-specific quality cuts rather than a single
blanket threshold, since the appropriate criterion depends on which
measurements an analysis uses. For continuum-derived quantities
(stellar mass, SFR, K-corrections), we recommend cutting on the
reduced $\chi^2$ of the continuum fit, \texttt{RCHI2\_CONT}~$<$
\todo{recommended threshold}. For individual emission lines, we
recommend selecting on line-amplitude signal-to-noise rather than
integrated-flux signal-to-noise---\todo{explain why amplitude S/N is
the more robust statistic (e.g., for broad or low-surface-brightness
lines)}---with a suggested threshold of \todo{X}$\sigma$;
\S\ref{sssec:kinematic}'s \texttt{DELTA\_LINECHI2} and
\texttt{DELTA\_LINENDOF} similarly quantify the significance of a
broad-line detection and can be used to build a custom broad-line
sample.

\todo{Is it worth adding a figure showing the correlation between
  amplitude and flux S/N?}

\subsection{Computational Performance}
\label{ssec:performance}

We run \fastspecfit{} on the NERSC Perlmutter supercomputer, whose
regular CPU nodes each provide two AMD EPYC 7763 (Milan) processors
with 64 cores each (128 physical cores, 256 logical threads per node
with SMT2). Fitting the full DR2 sample---38,245,183 objects in
52,325 HEALPix files---the code processes a typical DESI spectrum in
a median 4.8~s (mean 6.0~s) of single-core time, exploiting a hybrid
parallelism model: we use MPI to distribute independent HEALPix files
across ranks, and Python multiprocessing to parallelize the per-object
fits across all cores available to each rank. Per-object runtime
varies systematically with target class and observing
program---a mean 4.0~s/object/core for main-survey dark-time targets
(LRGs, ELGs, QSOs) versus 8.1~s/object/core for main-survey
bright-time targets (BGS, which typically require fitting more
emission lines)---and with the size of the HEALPix file being
processed, since \fastspecfit{} fits each file as a single in-memory
unit on one rank. Processing the complete DR2 catalog requires
approximately 850~node-hours on Perlmutter.

Because \fastspec{} fits tens of millions of spectra, we optimize the
per-object inner loops aggressively. We precompute and cache reused
arrays throughout---template FFTs for velocity-dispersion convolution
(\S\ref{sssec:vdisp}) and per-object redshift/IGM prefactors
(\S\ref{sssec:fitting}) chief among them---to avoid redundant
computation across the Monte Carlo realizations
(\S\ref{ssec:uncertainties}) and, for \vdisp, across the
$\chi^2$-scan grid. We implement the numerically intensive
kernels---dust attenuation, IGM transmission \citep{inoue14a},
resolution-matrix convolution, and the emission-line forward model
and its Jacobian---using Numba \citep[v\,\todo{X.Y};][]{lam15a}
just-in-time (JIT) compilation to native machine code. We represent
the emission-line Jacobian as an implicit matrix product (a
\texttt{scipy.sparse.linalg.LinearOperator}) rather than a dense
array, exploiting the fact that each truncated Gaussian line profile
($\pm5\sigma$) touches only a small, contiguous range of pixels; this
keeps both the memory footprint and the per-iteration cost of the
nonlinear least-squares solve far below what a dense Jacobian would
require. Because every JIT-compiled kernel uses Numba's persistent
on-disk compilation cache, we pay the one-time compilation cost once
per machine rather than once per process---a meaningful savings given
that our production run alone launches tens of thousands of
independent MPI ranks. \todo{Benchmark: X ms/spectrum, Y$\times$
  faster than NumPy.}

\subsection{Validation}
\label{ssec:validation}

\begin{figure*}
\includegraphics[width=6.5in]{compare-mstar-external.pdf}
\caption{Comparison of \fastspecfit{} stellar masses with three
  independent catalogs, organized by target class. Each row compares
  \fastspecfit{} $\logmstar$ (x-axis; $h=1$, Chabrier IMF) against
  masses from (top) \citet{zou24b} using CIGALE SED fitting of DR2
  photometry, (middle) \citet{siudek25a} using CIGALE SED fitting with
  AGN contributions fit to DR1 spectra and photometry, and (bottom)
  \citet{salim16a,salim18a} using CIGALE SED fitting of
  \textit{GALEX}$+$SDSS$+$\textit{WISE} photometry (GSWLC-X2). Columns
  show BGS, LRG, and ELG subsamples; the GSWLC-X2 comparison is
  restricted to BGS targets owing to negligible overlap with the other
  classes. The logarithmically-scaled two-dimensional number density
  is shown as a color map, with contours enclosing 50\%, 75\%, 95\%,
  and 99.5\% of the joint distribution; individual galaxies outside
  the outermost contour are shown as markers. Each panel lists the
  matched sample size $N$, the median mass offset $\Delta_{\rm med}$
  (external minus \fastspecfit{}) in dex, and the NMAD scatter. The
  dashed line is the 1:1 relation.\label{fig:compare-mstar-external}}
\end{figure*}

We validate the \fastspecfit{} measurements by comparing them against
independent codes and catalogs applied to overlapping subsets of the
DR2 sample: SED-fitting-based stellar masses and SFRs from three
external catalogs and from the COSMOS2020 photometric catalog,
stellar velocity dispersions from \texttt{pPXF}, and emission-line
fluxes from \texttt{pPXF}, the DESI \texttt{emlinefit} afterburner,
and GSWLC. For each comparison we select targets with reliable
measurements in both catalogs and quantify systematic offsets and
scatter as a function of signal-to-noise ratio, redshift, and target
class. \todo{Add summary of overall agreement and key findings.}

\begin{figure*}
\includegraphics[width=6.5in]{compare-sfr-external.pdf}
\caption{Write me\label{fig:compare-sfr-external}}
\end{figure*}

Figure~\ref{fig:compare-mstar-external} compares \fastspecfit{}
stellar masses against three independent CIGALE-based SED-fitting
catalogs. \todo{Describe agreement; highlight systematic offsets from
non-parametric vs. parametric SFH choices.} Figure~\ref{fig:compare-sfr-external}
makes the analogous comparison for star-formation rates. \todo{Describe
agreement.}

%Figure~\ref{fig:cosmos2020} compares \fastspecfit{} stellar masses
%against photometric stellar masses from the COSMOS2020 catalog
%\citep{weaver22a} for DESI SV3 targets in the COSMOS field, restricted
%to a photometric-redshift quality cut of \todo{specify cut}.
%\todo{Describe agreement; highlight systematic offsets from
%non-parametric vs. parametric SFH choices.}
%
%\begin{figure*}
%\includegraphics[width=6.5in]{compare-cosmos2020.pdf}
%\caption{Comparison of \fastspecfit{} stellar masses ($h=1$, Chabrier
%  IMF) with the COSMOS2020 photometric catalog \citep{weaver22a},
%  matched by sky position to DESI SV3 targets.  \textit{Top:}
%  COSMOS2020 $\logmstar$ versus \fastspecfit{} $\logmstar$; the
%  logarithmically-scaled two-dimensional number density is shown as a
%  color map with contours enclosing 50\%, 75\%, 95\%, and 99.5\% of
%  the joint distribution, and the solid line is the 1:1 relation.  The
%  panel lists the matched sample size $N$, the median offset
%  $\Delta_{\rm med}$ (COSMOS2020 minus \fastspecfit{}) in dex, and the
%  NMAD scatter.  \textit{Bottom:} Residual $\Delta\logmstar$ as a
%  function of redshift; the solid curve is the running median and the
%  dashed curves enclose the interquartile
%  range.\label{fig:cosmos2020}}
%\end{figure*}

Figure~\ref{fig:vdisp} compares \fastspecfit{} $\vdisp$ measurements
against those from \texttt{pPXF} \citep{cappellari23a} for a matched
sample of BGS galaxies with high signal-to-noise stellar
continua. \todo{Quantify agreement: rms, bias.}

\begin{figure*}
\includegraphics[width=6.5in]{compare-vdisp.pdf}
\caption{Comparison of \fastspecfit{} stellar velocity dispersions with
  those from \texttt{pPXF} \citep{cappellari23a} for a matched BGS
  sample.\label{fig:vdisp}}
\end{figure*}

We compare emission-line fluxes against \texttt{pPXF}
\citep{cappellari23a}, the DESI \texttt{emlinefit} afterburner
\citep{raichoor23a}, and the GSWLC catalog of \citet{salim16a} for
targets in the SDSS footprint. \todo{Describe agreement and key
  diagnostics compared (H$\alpha$, [\ion{O}{2}]); no figure currently
  implemented for this comparison in code/build-figures.py---confirm
  whether one should be built.}

%% ============================================================
\section{Discussion}
\label{sec:discussion}
%% ============================================================

The \fastspecfit{} DR2 catalog provides simultaneous measurements of
stellar continuum, emission-line, and broadband SED properties for
approximately 38~million extragalactic targets spanning four decades
in stellar mass and nearly two decades in redshift. In
\S\ref{ssec:reach} we illustrate the scientific reach of these
measurements with two representative figures: the distribution of
stellar mass with redshift across the three primary galaxy target
classes, and the optical emission-line diagnostics that
\fastspecfit{} uses to identify and characterize AGN. We then discuss
planned improvements to the code and future VACs
(\S\ref{ssec:improvements}).

\subsection{Galaxy Demographics and Diagnostic Power}
\label{ssec:reach}

\begin{figure*}
\includegraphics[width=6.5in]{figures/mstar-redshift.pdf}
\caption{Stellar mass versus redshift for BGS (left), LRG (center),
  and ELG (right) targets observed during the DESI SV3 bright and dark
  programs. The color scale and contours in each panel show the
  two-dimensional number density of galaxies passing per-class
  redshift quality cuts (see Section~X); contours enclose 50, 75, 95,
  and 99.5\% of the sample, and individual galaxies outside the
  outermost contour are plotted as points.
  %The number of galaxies in each panel is indicated.
  \label{fig:mstar_z}}
\end{figure*}

Figure~\ref{fig:mstar_z} shows the stellar mass--redshift distribution
for BGS, LRG, and ELG targets in the DR2 main survey. The three target
classes trace complementary regions of this plane: BGS spans the
widest mass range at $z \lesssim 0.6$, capturing galaxies from well
below the knee of the stellar mass function to the most massive
systems in the local universe; LRGs select the most massive, passively
evolving galaxies over $0.4 \lesssim z \lesssim 1.1$; and ELGs
preferentially identify low-to-intermediate mass star-forming systems
at $0.6 \lesssim z \lesssim 1.6$. Together these three samples provide
near-continuous coverage of the galaxy population from $z \approx 0$
to $z \approx 1.6$, making the DR2 catalog a uniquely comprehensive
resource for studies of galaxy evolution over the past $\sim$10~Gyr.

Figure~\ref{fig:bpt} shows the \citet{baldwin81a} BPT
diagram---\oiiilam/\hb{} versus \niilam/\ha---for BGS targets at $z
\lesssim 0.4$ with well-detected lines in all four diagnostics. The
demarcation curves of \citet{kewley01a} and \citet{kauffmann03c}
separate the star-forming sequence, composite systems, and AGN hosts;
the density distribution of the DR2 BGS sample across these regions
encodes the mix of nuclear activity in the low-redshift galaxy
population. \fastspecfit{}'s joint spectrophotometric fit enables
construction of this diagram for \todo{X}~million BGS galaxies in a
single, homogeneous analysis---a nearly tenfold increase over the SDSS
spectroscopic sample \todo{(cite)}. \todo{Quantify AGN, composite, and
  star-forming fractions; note redshift or mass dependence.}

%Where $P_1 = 0.63,\log[\ion{N}{2}]/\ha + 0.51,\log[\ion{S}{2}]/\ha +0.59,\log[\ion{O}{3}]/\hb$
and
%$P_3 = -0.46,\log[\ion{N}{2}]/\ha - 0.37,\log[\ion{S}{2}]/\ha + 0.81,\log[\ion{O}{3}]/\hb$

Be sure to cite \citet{ji20a} and \citet{jin21a}.

\begin{figure*}
\includegraphics[width=6.5in]{figures/bpt-agn.pdf}
\caption{AGN diagnostics for BGS galaxies in the DESI SV3 bright
  program. Both panels show the same sample of $N = {N}$ galaxies with
  signal-to-noise ratio $>3$ detected in all six diagnostic emission
  lines ([\ion{N}{2}] $\lambda6584$, \ha{}, [\ion{S}{2}]
  $\lambda\lambda6716,6731$, [\ion{O}{3}] $\lambda5007$, and
  \hb). The background color map and contours show the number
  density of objects. (Left) Classical BPT diagram
  \citep{baldwin81a}. The dashed curve shows the empirical
  star-forming/AGN demarcation of \citet{kauffmann03b}; galaxies above
  this line are AGN candidates. (Right) The $P_1$--$P_3$ projection of
  \citet{ji20a}, which separates star-forming galaxies (low $P_1$)
  from AGN (high $P_1$). The top axis shows the AGN fraction
  $f_\mathrm{AGN}$ mapped from $P_1$ following
  \citet{jin21a}. \label{fig:bpt}}
\end{figure*}

\subsection{Planned Improvements}
\label{ssec:improvements}

Several improvements to \fastspecfit{} are planned for future data
releases. On the modeling side, we plan to replace the current
solar-metallicity-only SPS templates (\S\ref{ssec:templates}) with
semi-empirical, non-negative matrix factorization (NMF) templates
trained directly on DESI spectra and spanning a range of stellar
metallicities; to adopt asymmetric emission-line profiles for lines
such as [\ion{O}{3}]~$\lambda5007$ that are known to exhibit
blue-wing outflow signatures; to forward-model Galactic foreground
reddening rather than correcting for it up front; and to improve the
aperture corrections that place fiber spectra on the total-flux scale
of the broadband photometry. We also plan to resolve a known
inconsistency between the emission lines fitted as part of the full
continuum model and those measured directly from the DESI spectrum,
including their line ratios \todo{cite Ashod's paper as a possible way
forward}.

On the fitting side, we plan to improve the treatment of QSOs and
broad-line AGN, including a dedicated \texttt{fastqso} fitting mode,
and to develop a new Bayesian fitting mode as an alternative to the
current maximum-likelihood approach (\S\ref{ssec:uncertainties}).
Finally, to keep pace with the growing scale of DESI and other
surveys, we plan to add GPU acceleration to the most computationally
intensive kernels (\S\ref{ssec:performance}) and to extend
\fastspecfit{} to the complete SDSS spectroscopic dataset.


%% ============================================================
\section{Summary}
\label{sec:summary}
%% ============================================================

\fastspecfit{} is a fast, flexible, and physically motivated code for
modeling DESI stellar continuum and emission-line spectra jointly with
broadband photometry. By convolving stellar-population and
emission-line templates with each fiber's native resolution matrix
(\S\S\ref{ssec:continuum}-\ref{ssec:emlines}) rather than resampling
the data, \fastspecfit{} preserves DESI's uncorrelated pixel
statistics while fitting a typical spectrum in a fraction of a second
on a single CPU core. Its non-parametric star-formation history,
energy-conserving dust model, and independently fitted narrow and
broad emission-line kinematics let the data---rather than a
pre-assumed galaxy type---determine the best-fitting model, so the
same machinery applies uniformly across the full diversity of DESI
extragalactic targets, from quiescent LRGs to star-forming ELGs to
broad-line QSOs (\S\ref{sec:casestudy}). Its speed and hybrid
MPI/multiprocessing design make it practical to estimate parameter
uncertainties from repeated Monte Carlo fits (\S\ref{ssec:uncertainties})
and to scale from a single laptop process to thousands of
supercomputer cores (\S\ref{ssec:performance}).

We use \fastspecfit{} to produce the DR2 VAC (Loa~v2.0), the primary
subject of this paper: approximately 38~million extragalactic targets
spanning the DESI commissioning, survey-validation, special-program,
and five-year main-survey campaigns (\S\ref{ssec:catalog}), including
corrected redshifts for roughly 165,000 QSOs. Processing the complete
catalog requires approximately 850~node-hours on the NERSC Perlmutter
supercomputer (\S\ref{ssec:performance}). We validate the resulting
stellar masses, star-formation rates, velocity dispersions, and
emission-line fluxes against independent codes and catalogs
(\S\ref{ssec:validation}), and we illustrate the catalog's scientific
reach with the stellar mass--redshift distribution of the three
primary galaxy target classes and with optical emission-line
diagnostics of nuclear activity (\S\ref{ssec:reach}). Companion
catalogs for the Early Data Release and Data Release~1, produced with
earlier template and \fastspecfit{} versions, are described in
Appendix~\ref{sec:appendix_vacs}.

Several improvements are planned for future \fastspecfit{} releases,
spanning the template library, the emission-line and dust models, new
fitting modes for QSOs and broad-line AGN, GPU acceleration, and an
extension to the complete SDSS spectroscopic dataset
(\S\ref{ssec:improvements}). The code and its measurements are already
enabling a broad range of extragalactic science using the public EDR
and DR1 VACs (\S\ref{sec:intro}), and a growing number of studies are
already drawing on the DR2 catalog described here, including
direct-method chemical abundance measurements for star-forming
galaxies \citep{scholte26a} and a search for the most metal-poor
star-forming galaxies yet found (J.~Moustakas et~al., in preparation);
we expect this number to grow substantially once the DR2 VAC is
publicly released in 2027.

The code is flexible! E.g., line-list, parameter constraints.

%% ============================================================
\begin{acknowledgements}
\JM{Need to thank Ben Johson. Should probably also add the Legacy
  Surveys acknowledgment.}

JM gratefully acknowledges funding support for this work from the
U.S. Department of Energy, Office of Science, Office of High Energy
Physics under Award Number DE-SC0020086.

This research is supported by the Director, Office of Science, Office of
High Energy Physics of the U.S. Department of Energy under Contract
No. DE-AC02-05CH11231, and by the National Energy Research Scientific
Computing Center, a DOE Office of Science User Facility under the same
contract; additional support for DESI is provided by the U.S. National
Science Foundation, Division of Astronomical Sciences under Contract
No. AST-0950945 to the NSF's National Optical-Infrared Astronomy Research
Laboratory; the Science and Technologies Facilities Council of the United
Kingdom; the Gordon and Betty Moore Foundation; the Heising-Simons
Foundation; the French Alternative Energies and Atomic Energy Commission
(CEA); the National Council of Science and Technology of Mexico; the
Ministry of Economy of Spain; and by the DESI Member Institutions.

The authors are honored to be permitted to conduct scientific research on
Iolkam Du'ag (Kitt Peak), a mountain with particular significance to the
Tohono O'odham Nation.
\end{acknowledgements}

\software{
  \texttt{Astropy} \citep{astropy-collaboration22a},
  \texttt{corner.py} \citep{corner},
  \texttt{healpy} \citep{zonca19a},
  \texttt{NumPy} \citep{harris20a},
  \texttt{numba} \citep{harris20a},
  \texttt{Matplotlib} \citep{hunter07a},
  \texttt{SciPy} \citep{virtanen20a},
  \texttt{seaborn} \citep{waskom21a}
}


%% ============================================================
\begin{appendix}

\section{DR2 VAC Data Model}
\label{sec:appendix_datamodel}

\todo{Possibly document the adaptive MORELINES HDU.}

We locate the per-HEALPix files, relative to the top-level VAC directory
\texttt{VACDIR}, at

\medskip
\noindent\texttt{VACDIR/healpix/SURVEY/PROGRAM/HPIXGROUP/HEALPIX/}\\
\texttt{\quad \{fastspec,fastphot\}-SURVEY-PROGRAM-HEALPIX.fits\{.gz\}}
\medskip

\noindent where \texttt{SURVEY}, \texttt{PROGRAM}, \texttt{HPIXGROUP}, and
\texttt{HEALPIX} follow the same conventions as the DESI data products. We
gzip-compress the \texttt{fastspec} HEALPix catalogs because they include a
\texttt{MODELS} extension containing the best-fitting model spectra (see
below); we do not compress the \texttt{fastphot} catalogs because they carry
no \texttt{MODELS} extension. We omit the \texttt{MODELS} extension from the
merged \texttt{fastspec} catalogs described in \S\ref{ssec:catalog} to keep
individual file sizes manageable. For large survey--program
combinations---specifically \texttt{main-bright} and \texttt{main-dark}---we
further subdivide the merged catalogs by \texttt{nside}~=~1 HEALPixel into 12
files following the naming pattern
\texttt{\{fastspec,fastphot\}-SURVEY-PROGRAM-nside1-hpNN.fits}, where
\texttt{NN} runs from \texttt{00} to \texttt{11}.

Each per-HEALPix \texttt{fastspec} file contains five extensions:
\begin{itemize}
  \item \texttt{HDU0} (\texttt{PRIMARY}): empty image with processing
    metadata keywords, including the software version, Monte Carlo
    random seed, and boolean fitting flags.
  \item \texttt{HDU1} (\texttt{METADATA}): target metadata mirroring
    the DESI redshift catalog, including sky coordinates, the input
    redshift and zwarning bitmask, targeting bits, and extinction-corrected
    Legacy Survey photometry.
  \item \texttt{HDU2} (\texttt{SPECPHOT}): derived spectrophotometric
    quantities, including SFH template coefficients, stellar-population
    parameters ($\tau_V$, $\vdisp$, age, metallicity, $\logmstar$, SFR),
    $D_n(4000)$, synthesized photometry, K-corrections, absolute
    magnitudes, monochromatic luminosities, and continuum flux densities
    at key emission-line centers.
  \item \texttt{HDU3} (\texttt{FASTSPEC}): emission-line fitting results,
    including per-camera median S/N and smooth continuum corrections,
    aperture correction factors, doublet line ratios, and per-line
    measurements (amplitude, flux, equivalent width, velocity shift,
    line width, and upper limits) for all modeled emission lines.
  \item \texttt{HDU4} (\texttt{MODELS}): 3D image array
    [$N_\lambda \times 3 \times N_\mathrm{obj}$] with the best-fitting
    model spectra on a common linear wavelength grid (3600--9824\,\AA,
    $0.8$\,\AA\,pixel$^{-1}$, vacuum wavelengths). The three planes
    store the full spectrophotometric model, the smooth continuum
    correction alone, and the emission-line model alone, respectively.
\end{itemize}

Table~\ref{table:datamodel} describes the key output columns in the
\texttt{SPECPHOT} and \texttt{FASTSPEC} extensions. All integrated
line fluxes are in units of \flux, flux densities in \fluxdensity, and
uncertainties are stored as inverse variances (a value of zero indicates
the parameter was not measured or the uncertainty is undefined).
A machine-readable description of the complete data model is maintained
in the \fastspecfit{} online
documentation.\footnote{\url{https://fastspecfit.readthedocs.io/en/latest/fastspec.html}}

\startlongtable
\begin{deluxetable*}{lll}
\tablecaption{FastSpecFit Output Data Model\label{table:datamodel}}
\tablewidth{0pt}
\tabletypesize{\scriptsize}
\tablehead{
  \colhead{Column} &
  \colhead{Units} &
  \colhead{Description}
}
\startdata
%% ------------------------------------------------------------------
\multicolumn{3}{c}{\textbf{METADATA Extension (HDU1)}} \\
%% ------------------------------------------------------------------
\texttt{TARGETID}                      & \nodata              & Unique DESI target identifier. \\
\texttt{RA}, \texttt{DEC}              & deg                  & Right ascension and declination. \\
\texttt{Z}                             & \nodata              & Redshift (Redrock, or QuasarNet for QSO targets). \\
\texttt{ZWARN}                         & \nodata              & Redrock zwarning bitmask. \\
\texttt{SPECTYPE}                      & \nodata              & Redrock spectral classification (\texttt{GALAXY}, \texttt{QSO}, \texttt{STAR}). \\
\texttt{FLUX\_[G/R/Z/W1/W2/W3/W4]}    & nMgy                 & Total photometric flux in each band, corrected for Galactic extinction. \\
\texttt{FLUX\_IVAR\_[G/R/Z/W1--W4]}   & nMgy$^{-2}$          & Inverse variance of \texttt{FLUX\_*}. \\
\texttt{FIBERFLUX\_[G/R/Z]}            & nMgy                 & Fiber flux corrected for Galactic extinction. \\
\texttt{EBV}                           & mag                  & Milky Way foreground $E(B-V)$ color excess. \\
\texttt{MW\_TRANSMISSION\_[G--W4]}     & \nodata              & Milky Way dust transmission factor $[0,1]$ per band. \\
%% ------------------------------------------------------------------
\multicolumn{3}{c}{\textbf{SPECPHOT Extension (HDU2)}} \\
%% ------------------------------------------------------------------
\multicolumn{3}{l}{\textit{Fit quality}} \\
\texttt{RCHI2}                         & \nodata              & Reduced $\chi^2$ of the full spectrophotometric fit. \\
\texttt{RCHI2\_CONT}                   & \nodata              & Reduced $\chi^2$ of the stellar continuum fit. \\
\texttt{RCHI2\_LINE}                   & \nodata              & Reduced $\chi^2$ of the emission-line fit. \\
\texttt{RCHI2\_PHOT}                   & \nodata              & Reduced $\chi^2$ of the photometric fit. \\
\multicolumn{3}{l}{\textit{Stellar continuum parameters}} \\
\texttt{COEFF[5]}                      & $\mathcal{F}_0\,M_\odot^{-1}$\tablenotemark{a} & Non-negative SFH template coefficients for the five lookback-time bins. \\
\texttt{VDISP} / \texttt{VDISP\_IVAR} & \kms                 & Stellar velocity dispersion and its inverse variance. \\
\texttt{TAUV} / \texttt{TAUV\_IVAR}   & \nodata              & $V$-band dust optical depth ($\tau_V$) and its inverse variance. \\
\texttt{AGE} / \texttt{AGE\_IVAR}     & Gyr                  & Light-weighted stellar age and its inverse variance. \\
\texttt{ZZSUN} / \texttt{ZZSUN\_IVAR} & \nodata              & Logarithmic stellar metallicity relative to solar and its inverse variance. \\
\texttt{LOGMSTAR} / \texttt{LOGMSTAR\_IVAR} & $\log(M_\odot)$ & Stellar mass ($h=1$, Chabrier IMF; eq.~\ref{eq:zfactor}) and its inverse variance. \\
\texttt{SFR} / \texttt{SFR\_IVAR}     & $M_\odot$\,yr$^{-1}$ & Instantaneous SFR (from youngest SFH bin) and its inverse variance. \\
\multicolumn{3}{l}{\textit{Spectral indices}} \\
\texttt{DN4000} / \texttt{DN4000\_IVAR}  & \nodata            & $D_n(4000)$ \citep{balogh99a} from the emission-line subtracted spectrum and its inverse variance. \\
\texttt{DN4000\_OBS}                   & \nodata              & $D_n(4000)$ measured from the observed spectrum (without emission-line subtraction). \\
\texttt{DN4000\_MODEL} / \texttt{DN4000\_MODEL\_IVAR} & \nodata & $D_n(4000)$ from the best-fitting continuum model and its inverse variance. \\
\multicolumn{3}{l}{\textit{Synthesized photometry}} \\
\texttt{FLUX\_SYNTH\_[G/R/Z]}          & nMgy                 & $g$-, $r$-, $z$-band flux synthesized from the observed spectrum. \\
\texttt{FLUX\_SYNTH\_SPECMODEL\_[G/R/Z]}  & nMgy             & $g$-, $r$-, $z$-band flux synthesized from the best-fitting spectroscopic model. \\
\texttt{FLUX\_SYNTH\_PHOTMODEL\_[G/R/Z/W1--W4]} & nMgy       & $grz$+W1--W4 flux synthesized from the best-fitting photometric model. \\
\multicolumn{3}{l}{\textit{Absolute magnitudes and K-corrections}\tablenotemark{b}} \\
\texttt{ABSMAG10\_DECAM\_[G/R/Z]}      & mag                  & DECam $g$-, $r$-, $z$-band absolute magnitude band-shifted to $z=1.0$. \\
\texttt{ABSMAG00\_[U/B/V]}             & mag                  & Johnson/Cousins $U$-, $B$-, $V$-band absolute magnitude band-shifted to $z=0$. \\
\texttt{ABSMAG00\_TWOMASS\_J}          & mag                  & 2MASS $J$-band absolute magnitude band-shifted to $z=0$. \\
\texttt{ABSMAG01\_SDSS\_[U/G/R/I/Z]}   & mag                  & SDSS $ugriz$ absolute magnitude band-shifted to $z=0.1$. \\
\texttt{ABSMAG01\_W1}                  & mag                  & WISE $W1$ absolute magnitude band-shifted to $z=0.1$. \\
\texttt{ABSMAG*\_IVAR}                 & mag$^{-2}$           & Inverse variance of the corresponding absolute magnitude. \\
\texttt{ABSMAG*\_SYNTH}                & mag                  & Synthesized absolute magnitude from the best-fitting SED model. \\
\texttt{KCORR10\_DECAM\_[G/R/Z]}       & mag                  & K-correction for DECam $grz$ absolute magnitudes at $z=1.0$. \\
\texttt{KCORR00\_[U/B/V]}              & mag                  & K-correction for Johnson/Cousins $UBV$ at $z=0$. \\
\texttt{KCORR01\_SDSS\_[U/G/R/I/Z]}    & mag                  & K-correction for SDSS $ugriz$ at $z=0.1$. \\
\texttt{KCORR01\_W1}                   & mag                  & K-correction for WISE $W1$ at $z=0.1$. \\
\multicolumn{3}{l}{\textit{Monochromatic luminosities}} \\
\texttt{LOGLNU\_1500} / \texttt{LOGLNU\_1500\_IVAR} & $10^{-28}$\,erg\,s$^{-1}$\,Hz$^{-1}$ & Monochromatic luminosity at rest-frame 1500\,\AA\ and its inverse variance. \\
\texttt{LOGLNU\_2800} / \texttt{LOGLNU\_2800\_IVAR} & $10^{-28}$\,erg\,s$^{-1}$\,Hz$^{-1}$ & Monochromatic luminosity at rest-frame 2800\,\AA\ and its inverse variance. \\
\texttt{LOGL\_[1450/1700/3000/5100]}   & $10^{10}\,L_\odot$   & Integrated luminosity at 1450, 1700, 3000, and 5100\,\AA\ in the rest-frame. \\
\texttt{LOGL\_[1450/1700/3000/5100]\_IVAR} & $10^{-20}\,L_\odot^{-2}$ & Inverse variance of the corresponding integrated luminosity. \\
\multicolumn{3}{l}{\textit{Continuum flux densities at key emission-line centers}} \\
\texttt{FLYA\_1215\_CONT} / \texttt{..\_IVAR} & \fluxdensity & Continuum flux density at Ly$\alpha$\,$\lambda$1215.67 and its inverse variance. \\
\texttt{FOII\_3727\_CONT} / \texttt{..\_IVAR} & \fluxdensity & Continuum flux density at {[\ion{O}{2}]}\,$\lambda$3728 and its inverse variance. \\
\texttt{FHBETA\_CONT} / \texttt{..\_IVAR}     & \fluxdensity & Continuum flux density at H$\beta$\,$\lambda$4862 and its inverse variance. \\
\texttt{FOIII\_5007\_CONT} / \texttt{..\_IVAR} & \fluxdensity & Continuum flux density at {[\ion{O}{3}]}\,$\lambda$5007 and its inverse variance. \\
\texttt{FHALPHA\_CONT} / \texttt{..\_IVAR}    & \fluxdensity & Continuum flux density at H$\alpha$\,$\lambda$6564 and its inverse variance. \\
%% ------------------------------------------------------------------
\multicolumn{3}{c}{\textbf{FASTSPEC Extension (HDU3)}} \\
%% ------------------------------------------------------------------
\multicolumn{3}{l}{\textit{Spectral diagnostics}} \\
\texttt{SNR\_[B/R/Z]}                  & \nodata              & Median signal-to-noise ratio per pixel in the $B$, $R$, $Z$ cameras. \\
\texttt{SMOOTHCORR\_[B/R/Z]}           & \%                   & Median smooth continuum correction relative to data in each camera (\S\ref{ssec:smooth}). \\
\texttt{APERCORR}                      & \nodata              & Median aperture correction factor (\S\ref{ssec:fitting}). \\
\texttt{APERCORR\_[G/R/Z]}             & \nodata              & Per-band aperture correction factor. \\
\multicolumn{3}{l}{\textit{Broad-line model selection}} \\
\texttt{DELTA\_LINECHI2}               & \nodata              & $\Delta\chi^2$ between narrow-only and narrow+broad emission-line models. \\
\texttt{DELTA\_LINENDOF}               & \nodata              & Corresponding difference in degrees of freedom. \\
\multicolumn{3}{l}{\textit{Doublet line ratios}} \\
\texttt{MGII\_DOUBLET\_RATIO} / \texttt{..\_IVAR}   & \nodata & \ion{Mg}{2} 2796/2803 flux ratio and its inverse variance. \\
\texttt{OII\_DOUBLET\_RATIO} / \texttt{..\_IVAR}    & \nodata & {[\ion{O}{2}]} 3726/3729 flux ratio and its inverse variance. \\
\texttt{OIII\_DOUBLET\_RATIO} / \texttt{..\_IVAR}   & \nodata & {[\ion{O}{3}]} 5007/4959 flux ratio and its inverse variance. \\
\texttt{SII\_DOUBLET\_RATIO} / \texttt{..\_IVAR}    & \nodata & {[\ion{S}{2}]} 6731/6716 flux ratio and its inverse variance. \\
\texttt{NII\_DOUBLET\_RATIO} / \texttt{..\_IVAR}    & \nodata & {[\ion{N}{2}]} 6584/6548 flux ratio and its inverse variance. \\
\texttt{OIIRED\_DOUBLET\_RATIO} / \texttt{..\_IVAR} & \nodata & {[\ion{O}{2}]} 7320/7330 flux ratio and its inverse variance. \\
\multicolumn{3}{l}{\textit{Emission-line measurements (\texttt{LINENAME}\tablenotemark{c})}} \\
\texttt{LINENAME\_MODELAMP}            & \fluxdensity         & Model line amplitude before convolution with the resolution matrix (eq.~\ref{eq:intrinsic}). \\
\texttt{LINENAME\_AMP} / \texttt{..\_IVAR} & \fluxdensity    & Observed emission-line amplitude and its inverse variance. \\
\texttt{LINENAME\_FLUX} / \texttt{..\_IVAR} & \flux          & Gaussian-integrated emission-line flux (eq.~\ref{eq:flux_phys}) and its inverse variance. \\
\texttt{LINENAME\_BOXFLUX} / \texttt{..\_IVAR} & \flux       & Boxcar-integrated emission-line flux and its inverse variance. \\
\texttt{LINENAME\_VSHIFT} / \texttt{..\_IVAR} & \kms         & Velocity shift relative to \texttt{Z} (eq.~\ref{eq:lstar}) and its inverse variance. \\
\texttt{LINENAME\_SIGMA} / \texttt{..\_IVAR}  & \kms         & Intrinsic Gaussian line width before resolution convolution and its inverse variance. \\
\texttt{LINENAME\_CONT} / \texttt{..\_IVAR}   & \fluxdensity & Continuum flux density at the line center and its inverse variance. \\
\texttt{LINENAME\_EW} / \texttt{..\_IVAR}     & \AA          & Rest-frame equivalent width and its inverse variance. \\
\texttt{LINENAME\_FLUX\_LIMIT}         & erg\,s$^{-1}$\,cm$^{-2}$ & $1\sigma$ upper limit on the emission-line flux. \\
\texttt{LINENAME\_CHI2}                & \nodata              & Chi-squared of the line fit. \\
\texttt{LINENAME\_NPIX}                & \nodata              & Number of pixels attributed to the emission line. \\
\multicolumn{3}{l}{\textit{Non-parametric line moments (\texttt{LINENAME2}\tablenotemark{d})}} \\
\texttt{LINENAME2\_MOMENT1} / \texttt{..\_IVAR} & \AA        & First moment (centroid) within $\pm5\sigma$ of the line center and its inverse variance. \\
\texttt{LINENAME2\_MOMENT2} / \texttt{..\_IVAR} & \AA$^2$   & Second moment (variance) relative to \texttt{MOMENT1} and its inverse variance. \\
\texttt{LINENAME2\_MOMENT3} / \texttt{..\_IVAR} & \AA$^3$   & Third moment (skewness) relative to \texttt{MOMENT1} and its inverse variance. \\
\enddata
\tablecomments{The \texttt{MODELS} extension (HDU4) contains a 3D image array [$N_\lambda \times 3 \times N_\mathrm{obj}$] with the best-fitting model spectra on a common linear wavelength grid (3600--9824\,\AA, $0.8$\,\AA\,pixel$^{-1}$). The three planes store the full spectrophotometric model, the smooth continuum correction alone, and the emission-line model alone, respectively. Fluxes are in units of $10^{-17}$\,erg\,s$^{-1}$\,cm$^{-2}$\,\AA$^{-1}$. Uncertainties for all parameters are stored as inverse variances; a value of zero indicates the parameter was not measured.}
\tablenotetext{a}{$\mathcal{F}_0 = 10^{-17}$\,erg\,s$^{-1}$\,cm$^{-2}$\,\AA$^{-1}$. Template coefficients have units of $\mathcal{F}_0\,M_\odot^{-1}$, so multiplying by the template normalization mass $\mathcal{M}_0 = 10^{10}\,M_\odot$ gives $\mathcal{F}_0$, consistent with the stellar mass in eq.~(\ref{eq:zfactor}).}
\tablenotetext{b}{Absolute magnitudes are derived from the observed-frame bandpass nearest in wavelength to the desired rest-frame band, provided $\mathrm{S/N} > 2$. If no photometry satisfies this criterion, the magnitude is synthesized from the best-fitting model and its inverse variance is set to zero.}
\tablenotetext{c}{\texttt{LINENAME} represents each of the \todo{N} emission lines modeled by \fastspecfit{}, including broad Balmer components; see Table~\ref{table:emlines}.}
\tablenotetext{d}{\texttt{LINENAME2} represents the subset of four lines for which non-parametric moments are computed: \ion{C}{4}\,$\lambda$1549, \ion{Mg}{2}\,$\lambda$2800, H$\beta$, and {[\ion{O}{3}]}\,$\lambda$5007.}
\end{deluxetable*}


\section{Pipeline Diagnostics}
\label{sec:appendix_diagnostics}

One of the dividends of fitting every DESI spectrum with a physically
motivated forward model is that the residuals carry information about
the quality of the underlying data reduction. In this appendix we use
\fastspecfit{} measurements to characterize two classes of residual
systematic error in the DR2 pipeline: wavelength-solution errors
revealed by per-camera emission-line velocity shifts
(\S\ref{ssec:waveresid}), and spectrophotometric calibration residuals
revealed by the difference between synthesized and observed broadband
photometry as a function of magnitude, redshift, and target class
(\S\ref{ssec:fluxresid}).

\subsection{Wavelength Calibration Residuals}
\label{ssec:waveresid}

\fastspecfit{} fits each of the three DESI cameras (B, R, Z) separately
and records the median emission-line velocity shift per camera for every
spectrum in which emission lines are detected. For a perfectly calibrated
wavelength solution, this per-camera shift should be consistent with zero
and with the global redshift measured by Redrock. Systematic offsets or
trends with wavelength, fiber position, or observing conditions indicate
residual errors in the arc-lamp or sky-line wavelength calibration.

\begin{figure*}
%\includegraphics[width=6.5in]{figures/wave-residuals.png}
\caption{Per-camera median emission-line velocity shifts as a function of
  wavelength and target class for a representative subset of the DR2
  sample. Each panel shows one DESI camera (B, R, Z); the dashed line
  marks zero offset. \todo{Generate figure.}\label{fig:wave_resid}}
\end{figure*}

Figure~\ref{fig:wave_resid} shows the distribution of per-camera velocity
shifts for emission-line targets across the full DR2 footprint as a
function of \todo{relevant variable: fiber number, focal-plane position,
observing date?}. \todo{Describe findings: are the residuals consistent
with zero? Any trends with target class, S/N, or observing conditions?}

\subsection{Spectrophotometric Calibration Residuals}
\label{ssec:fluxresid}

\fastspecfit{} simultaneously fits the DESI spectra and the Legacy Survey
broadband photometry (\S\ref{ssec:photometry}), so the difference between
the synthesized and observed photometry is a direct measure of residual
spectrophotometric calibration errors in the reduced spectra. We quantify
these residuals as $\Delta m = m_\mathrm{synth} - m_\mathrm{obs}$ in each
Legacy Survey band and examine their dependence on apparent magnitude,
redshift, and target class.

\begin{figure*}
%\includegraphics[width=6.5in]{figures/fluxcal-residuals.png}
\caption{Residual spectrophotometric calibration ($\Delta m =
  m_\mathrm{synth} - m_\mathrm{obs}$) as a function of apparent magnitude
  and redshift for BGS, LRG, and ELG targets. Each column corresponds to
  a Legacy Survey band ($g$, $r$, $z$); rows separate target class. Dashed
  lines mark zero offset; contours enclose 68 and 95 per cent of the
  sample. \todo{Generate figure.}\label{fig:fluxcal_resid}}
\end{figure*}

Figure~\ref{fig:fluxcal_resid} shows $\Delta m$ as a function of
\todo{magnitude and/or redshift} for the three primary extragalactic
target classes. \todo{Describe findings: systematic offsets? Magnitude
or color dependence? Variation across target classes?}

%% ============================================================
\section{Stellar Masses inferred using \fastphot}
\label{sec:appendix_fastphot}

\begin{figure*}
\includegraphics[width=6.5in]{figures/compare-mstar-fastphot.pdf}
\caption{Comparison of \fastspecfit{} stellar masses derived from
  spectrophotometric fitting (\fastspec) and from photometry-only
  fitting (\fastphot{}), organized by target class.  \textit{Top row:}
  \fastphot{} $\logmstar$ versus \fastspec{} $\logmstar$ ($h=1$,
  Chabrier IMF) for BGS (left), LRG (center), and ELG (right)
  subsamples.  The logarithmically-scaled two-dimensional number
  density is shown as a color map, with contours enclosing 50\%, 75\%,
  95\%, and 99.5\% of the joint distribution; individual galaxies
  outside the outermost contour are shown as markers.  Each panel
  lists the matched sample size $N$, the median mass offset
  $\Delta_{\rm med}$ (\fastphot{} minus \fastspec) in dex, and the
  NMAD scatter.  The dashed line is the 1:1 relation.  \textit{Bottom
    row:} Residual $\Delta\logmstar$ (\fastphot{} minus \fastspec)
  as a function of redshift for the same three subsamples.  The solid
  curve shows the running median and the shaded band encloses the
  interquartile range.  The dashed horizontal line marks zero offset.
  We attribute the systematic offset seen in BGS targets at low
  redshift to imperfect aperture corrections for the most extended,
  nearby galaxies, whose half-light radii subtend a significant
  fraction of the DESI fiber
  diameter.\label{fig:compare_mstar_fastphot}}
\end{figure*}

Figure~\ref{fig:compare_mstar_fastphot} compares the stellar masses we derive
using \fastspec{} and \fastphot{} for targets in common, illustrating
the consistency between the spectrophotometric and photometry-only
fitting modes.



%% ============================================================
\section{EDR and DR1 Value-Added Catalogs}
\label{sec:appendix_vacs}

Two known issues affect the v1.0 DR2 catalog. First, the
inverse-variance estimates for emission-line fluxes and equivalent
widths carry significant errors; users requiring precise flux
uncertainties should apply appropriate caution and consult the online
documentation for the latest status of corrections. Second, a small
fraction of spectra with strong emission lines suffer poor line fits
due to a clipping error in the internal flux uncertainty routine. We
fix both issues in \fastspecfit{}~v3.2.0, and we will release updated
catalogs in due course.

The EDR VAC (Fuji~v3.2) comprises approximately 1.4~million targets
processed with \fastspecfit{}~v2.5.0 and was publicly released in
January~2024. The DR1 VAC (Iron~v3.0) adds the first tranche of
main-survey data, processed with \fastspecfit{}~v3.1.5, and totals
approximately 18~million targets; it was publicly released in
June~2025. Table~\ref{table:sample-edr-dr1} summarizes the EDR and
DR1 target counts organized by survey and observing program, for
direct comparison with the DR2 breakdown in Table~\ref{table:sample}.

\begin{deluxetable*}{lrr}
\tablecaption{Summary of the EDR and DR1 \fastspecfit{} Value-Added Catalogs\label{table:sample-edr-dr1}}
\tablewidth{0pt}
\tablehead{
  \colhead{Program} &
  \colhead{EDR\tablenotemark{a}} &
  \colhead{DR1\tablenotemark{b}}
}
\startdata
\cutinhead{Main Survey}
backup &   \nodata &     15,163 \\
bright &   \nodata &  6,445,927 \\
dark   &   \nodata & 10,075,629 \\
total  &   \nodata & 16,536,719 \\
\cutinhead{Special Survey}
backup &   \nodata &        552 \\
bright &   \nodata &     43,261 \\
dark   &    35,649 &     14,954 \\
other  &   \nodata &     42,064 \\
total  &    35,649 &    100,831 \\
\cutinhead{Survey Validation\tablenotemark{c}}
backup &     5,354 &      4,960 \\
bright &   438,513 &    438,429 \\
dark   &   881,164 &    878,083 \\
other  &    36,923 &     36,874 \\
total  & 1,361,954 &  1,358,346 \\
\tableline
Total & 1,397,603 & 17,995,896 \\
\enddata
\tablenotetext{a}{Fuji v3.2; released 2024 January.}
\tablenotetext{b}{Iron v3.0; released 2025 June.}
\tablenotetext{c}{Sum of CMX, SV1, SV2, and SV3 surveys.}
\tablecomments{Target counts in the \fastspecfit{} EDR and DR1 VACs, organized
  by program. See Table~\ref{table:sample} for the corresponding DR2 counts.}
\end{deluxetable*}


Relative to DR2, the EDR and DR1 catalogs are produced with earlier
versions of the SPS template library and do not incorporate the
corrected QSO redshifts described in \S\ref{ssec:catalog} or the
flux-uncertainty and line-fitting fixes introduced in
\fastspecfit{}~v3.2.0. Users combining measurements across data
releases should therefore expect small, systematic differences in
stellar masses, star-formation rates, and emission-line properties,
particularly for AGN and QSO targets. We defer a full quantitative
comparison of stellar masses and emission-line measurements across
releases to future work.

\end{appendix}


\bibliographystyle{aasjournal}
\bibliography{refs}

\end{document}
